% 火星（综述）
% license CCBYSA3
% type Wiki

本文根据 CC-BY-SA 协议转载翻译自维基百科\href{https://en.wikipedia.org/wiki/Mars}{相关文章}。

\begin{figure}[ht]
\centering
\includegraphics[width=6cm]{./figures/b86f152bd3571454.png}
\caption{火星的真实颜色\(^{\text{a}}\)，由“希望号”轨道器拍摄。画面中央可见塔尔西斯山脉，左侧为奥林匹斯山，右侧则是水手谷。} \label{fig_Mars_1}
\end{figure}
火星是距离太阳的第四颗行星。由于其呈橙红色的外观，它也被称为“红色星球”\(^{\text{[22,23]}}\)。火星是一颗类似沙漠的岩质行星，大气稀薄，主要成分为二氧化碳（CO\(_2\)）。在平均地表高度处，其大气压仅为地球的几千分之一，大气温度范围为 −153 至 20 °C（−243 至 68 °F）\(^{\text{[24]}}\)，宇宙辐射水平很高。火星保留了一部分水，存在于地下以及稀薄大气中，可形成卷云、雾气、霜冻、较大的极地永久冻土区与冰帽（伴随季节性的 CO\(_2\) 雪），但其表面不存在液态水体。其表面重力约为地球的三分之一，为月球的两倍。直径为 6,779 公里（4,212 英里），约为地球的一半，是月球的两倍，其表面积则相当于地球全部陆地面积的总和。

细微尘埃遍布火星表面和大气，在低重力条件下，即便是稀薄大气的微弱风力也能将其扬起并扩散开来。火星地形大致呈现显著的南北分界，即“火星二元性结构”：北半球主要是平坦、低洼的平原，而南半球则由遍布陨石坑的高地组成。从地质学上看，火星处于相对活跃的状态，地下会发生“火震”，但其上也分布着许多巨型而已灭绝的火山（最高的是奥林匹斯山，高 21.9 公里或 13.6 英里），以及太阳系中最大峡谷之一的水手谷（长约 4,000 公里或 2,500 英里）。火星有两颗天然卫星，形状小而不规则：火卫一（Phobos）与火卫二（Deimos）。由于其轴倾角为 25 度，火星像地球（轴倾角 23.5 度）一样经历季节变化。一个火星太阳年相当于 1.88 个地球年（687 个地球日），一个火星太阳日（sol）则等于 24.6 小时。

火星与其他行星一同形成于大约 45 亿年前。在火星的诺亚纪（Noachian period，约 45 亿年至 35 亿年前）期间，其地表经历了陨石撞击、峡谷形成、侵蚀作用、可能存在的海洋，以及磁层的消失。赫斯珀里亚纪（Hesperian period，自约 35 亿年前开始，结束于约 33–29 亿年前）以广泛的火山活动和导致巨大泄流河道形成的洪水为主要特征。亚马逊纪（Amazonian period）从其后延续至今，是当前地质过程的主要影响因素。由于火星的地质历史，对于火星过去或当前是否存在生命的探索仍是活跃的科学研究领域，其中一些可能的生命迹象仍需进一步确认。

由于在地球夜空中可被肉眼看到、呈红色并缓慢移动，火星自古以来就被人类观测，并在不同文化中具有多样的象征意义。1963 年，人类首次发射前往火星的探测器“火星 1 号”（Mars 1），但途中失联。第一项成功飞越火星的探测任务发生在 1965 年，由“水手 4 号”（Mariner 4）完成。1971 年，“水手 9 号”（Mariner 9）进入火星轨道，成为人类历史上首个进入除月球、太阳或地球以外天体轨道的航天器；同年，“火星 2 号”成为首个在火星发生非受控撞击的探测器，“火星 3 号”则实现了首次成功着陆。自 1997 年以来，探测器在火星上持续运行。有时，超过十个探测器会同时在火星轨道或地表开展活动，这一数量超过地球以外任意行星。

火星常被视为未来载人探索任务的目标，尽管目前尚无正式计划付诸实施。
\subsection{自然历史}  
\subsubsection{形成} 
科学家推测，在太阳系形成过程中，火星是由于围绕太阳运行的原行星盘中的物质通过随机的失控聚积过程形成的。火星在太阳系中的位置使其具备许多独特的化学特征。相较于地球，沸点较低的元素，例如氯、磷和硫，在火星上要常见得多；这些元素很可能是在年轻太阳强烈的太阳风作用下被推移到更外侧区域\(^{\text{[25]}}\)。

晚期重轰炸  
在行星形成之后，内太阳系可能经历了所谓的“晚期重轰炸”。火星约有 60\% 的表面保留了那个时期的撞击记录\(^{\text{[26,27,28]}}\)，而其余许多区域的地表之下，可能也存在由同一时期事件形成的巨大撞击盆地。然而，最近的建模研究对晚期重轰炸的存在提出了质疑\(^{\text{[29]}}\)。证据显示，火星北半球可能存在一个巨大的撞击盆地，范围达 10,600 × 8,500 公里（6,600 × 5,300 英里），约为月球南极—艾特肯盆地面积的四倍；若得到确认，这将是迄今发现的最大撞击盆地\(^{\text{[30]}}\)。科学家推测该盆地形成于约 40 亿年前，当时火星被一个与冥王星大小相当的天体撞击，造成了火星南北半球的显著地形差异，并形成覆盖火星 40\% 表面的平滑的北方平原（Borealis basin）\(^{\text{[31,32]}}\)。

一项 2023 年的研究根据火卫二（Deimos）的轨道倾角提出证据，显示火星在 35 亿至 40 亿年前可能曾拥有一个环系统\(^{\text{[33]}}\)。该环系统可能由一颗质量约为火卫一（Phobos）20 倍的古老卫星解体形成；而火卫一可能正是这一原始环系统的残余物\(^{\text{[34,35]}}\)。
\begin{figure}[ht]
\centering
\includegraphics[width=14.25cm]{./figures/4b6ac302fa80c765.png}
\caption{} \label{fig_Mars_2}
\end{figure}
\begin{figure}[ht]
\centering
\includegraphics[width=14.25cm]{./figures/7861217b62cae434.png}
\caption{} \label{fig_Mars_3}
\end{figure}
火星的地质历史可以被划分为多个时期，但以下三个时期最为主要 \(^{\text{[36,37]}}\)。
\begin{itemize}
\item 诺亚纪（Noachian period）：约 45 亿年至 35 亿年前，是火星现存最古老地表的形成阶段。诺亚纪地表布满大量大型撞击坑。塔尔西斯隆起（Tharsis bulge）这一火山高地被认为形成于该时期，其末期可能经历了大规模液态水泛滥。该时期名称源自 Noachis Terra \(^{\text{[38]}}\)。
\item 赫斯珀纪（Hesperian period）：约 35 亿年至介于 33 亿年至 29 亿年前之间。赫斯珀纪的特征是广泛的熔岩平原的形成。名称源自 Hesperia Planum \(^{\text{[38]}}\)。
\item 亚马孙纪（Amazonian period）：约 33 亿年至 29 亿年前至今。亚马孙纪区域的陨石撞击坑数量较少，但地貌表现多样。奥林匹斯山（Olympus Mons）形成于该时期，同时火星其他地区也存在熔岩流活动。名称源自 Amazonis Planitia \(^{\text{[38]}}\)。
\end{itemize}
\subsubsection{近期地质活动}
火星上仍然存在地质活动。阿萨巴斯卡峡谷（Athabasca Valles）中有形成于约 2 亿年前的片状熔岩流。刻耳柏洛斯裂谷（Cerberus Fossae）这一断陷带中的水流活动发生于不足 2,000 万年前，显示出同样年轻的火山侵入事件 \(^{\text{[39]}}\)。火星勘测轨道器（Mars Reconnaissance Orbiter）已拍摄到雪崩的影像 \(^{\text{[40,41]}}\)。
\subsection{物理特征}
\begin{figure}[ht]
\centering
\includegraphics[width=8cm]{./figures/23fb842f59355847.png}
\caption{按比例绘制的火星及内太阳系中具有行星质量的天体。从左至右依次为：水星、金星、地球、月球、火星和谷神星。} \label{fig_Mars_4}
\end{figure}
火星的直径约为地球的一半，或约为月球的两倍，其表面积仅略小于地球全部陆地面积 \(^{\text{[2]}}\)。火星的密度低于地球，体积约为地球的 15\%，质量约为地球的 11\%，因此其表面重力约为地球的 38\%。火星是目前已知的唯一一颗沙漠行星，即表面类似于地球荒漠环境的类地岩质行星。火星表面呈红橙色，是由于铁(III)氧化物（纳米级 Fe\(_2\)O\(_3\)）以及铁(III)氧化物-氢氧化物矿物针铁矿（goethite）所致 \(^{\text{[42]}}\)。其颜色有时类似奶油太妃糖 \(^{\text{[43]}}\)；其他常见的地表颜色还包括金色、棕色、褐色以及因矿物组成不同而呈现的绿ish色调 \(^{\text{[43]}}\)。
\subsubsection{内部结构}
\begin{figure}[ht]
\centering
\includegraphics[width=8cm]{./figures/4694e492a3c0d60d.png}
\caption{火星的内部结构 \(^{\text{[44,45]}}\)} \label{fig_Mars_5}
\end{figure}
与地球类似，火星也分异为高密度的金属核以及其上覆盖的低密度岩质层 \(^{\text{[46][47]}}\)。最外层为地壳，平均厚度约为 42–56 公里（26–35 英里） \(^{\text{[48]}}\)，在伊西迪斯平原（Isidis Planitia）处最薄，为 6 公里（3.7 英里），在塔尔西斯高原（Tharsis plateau）南部最厚，可达 117 公里（73 英里） \(^{\text{[49]}}\)。相比之下，地球地壳平均厚度为 27.3 ± 4.8 公里 \(^{\text{[50]}}\)。火星地壳中最丰富的元素包括硅、氧、铁、镁、铝、钙和钾。火星已被确认具有地震活动 \(^{\text{[51]}}\)；在 2019 年，报道指出 InSight 着陆器已探测并记录超过 450 次火震及相关事件 \(^{\text{[52][53]}}\)。

在地壳之下是硅酸盐地幔，它形成了火星表面许多构造与火山特征。火星上地幔的上部为低速带，在此区段中地震波速度低于邻近深度层。地幔似乎在约 250 公里深度范围内保持刚性 \(^{\text{[45]}}\)，使得火星的岩石圈相较地球更为厚实。在此深度以下，地幔逐渐变得更具延展性，且地震波速度再次开始增加 \(^{\text{[54]}}\)。火星地幔似乎并不存在类似地球下地幔的热绝缘层；相反，在约 1050 公里深度以下，其矿物学特征与地球的过渡带相似 \(^{\text{[44]}}\)。在地幔底部存在一层约 150–180 公里厚的基性液态硅酸盐层 \(^{\text{[45][55]}}\)。火星地幔看似高度非均质，其中包含直径可达 4 公里的高密度碎块，可能是在约 45 亿年前的巨型撞击中被深度注入；来自八次火震的高频波在穿越这些局部区域时出现减速，模型显示这些非均质体是组成上独特的残余碎屑，由于火星缺乏板块构造且内部对流迟缓，无法实现完全均质化，因此得以保存 \(^{\text{[56][57]}}\)。

火星的铁镍核至少部分为熔融状态，并可能具有固体内核 \(^{\text{[58][59][60]}}\)。其半径约为火星半径的一半，大约为 1650–1675 公里，且富含硫、氧、碳和氢等轻元素 \(^{\text{[61][62]}}\)。其核心温度估计为 2000–2400 K \(^{\text{[63]}}\)，相比之下，地球固体内核的温度为 5400–6230 K。2025 年，基于 InSight 着陆器的数据，一组研究人员报告探测到一个半径为 613 公里（381 英里） ± 67 公里（42 英里）的固体内核 \(^{\text{[60]}}\)。
\subsubsection{表面地质}  
更多信息参见：火星风化层（Martian regolith）。火星是一颗类地行星，其表面由含硅和氧的矿物、金属以及其他通常构成岩石的元素组成。火星表面主要由拉斑玄武岩（tholeiitic basalt）构成 \(^{\text{[64]}}\)，但部分区域的含硅量高于典型玄武岩，可能类似地球的安山岩或硅质玻璃。低反照率区域显示斜长石长石的富集，而北部低反照率地区呈现较高含量的层状硅酸盐与高硅玻璃。南部高地的部分地区可检测到高钙辉石的存在。此外，还发现局地富集的赤铁矿与橄榄石 \(^{\text{[65]}}\)。火星大部分表面深覆以细粒的铁(III)氧化物尘埃 \(^{\text{[66]}}\)。

凤凰号着陆器（Phoenix lander）返回的数据表明火星土壤略呈碱性，且含有镁、钠、钾及氯等元素。这些养分同样存在于地球土壤中，并对植物生长必不可少 \(^{\text{[67]}}\)。着陆器的实验显示火星土壤的 pH 为 7.7，为弱碱性，并按重量含有 0.6\% 的高氯酸盐 \(^{\text{[68][69]}}\)，这种浓度对人类具有毒性 \(^{\text{[70][71]}}\)。

条纹（streaks）在火星上十分常见，并且在陨石坑、槽谷与峡谷的陡坡上频繁出现新的条纹。这些条纹初期颜色较深，随后随时间变浅。条纹可从极小区域开始，并扩展数百米，其路径会绕过巨石等障碍物。最普遍接受的假说认为，这些条纹是由明亮尘埃的雪崩或尘卷风作用暴露出的暗色下层土壤 \(^{\text{[72]}}\)。其他提出的解释包括与水有关的机制，甚至可能涉及生物生长的假说 \(^{\text{[73][74]}}\)。

火星表面的环境辐射水平平均为每日 0.64 毫西弗，显著低于往返火星途中每日 1.84 毫西弗（或每日 22 毫拉德）的辐射水平 \(^{\text{[75][76]}}\)。相比之下，地球近地轨道（低地轨道）中，例如空间站所在轨道的辐射约为每日 0.5 毫西弗 \(^{\text{[77]}}\)。赫拉斯平原（Hellas Planitia）具有火星表面最低的辐射水平，约每日 0.342 毫西弗，其位于哈德里亚卡斯山（Hadriacus Mons）西南方的熔岩管（lava tubes）中辐射水平可能低至每日 0.064 毫西弗 \(^{\text{[78]}}\)，与地球飞机飞行期间的辐射水平相当。
\subsubsection{磁学特征} 
虽然火星不存在有组织的全球磁场的证据 \(^{\text{[79]}}\)，但观测表明其地壳的部分区域已被磁化，这暗示其偶极磁场在过去曾发生过极性反转。磁化敏感矿物所记录的这种古地磁性与地球洋底发现的交替磁带类似。一项发表于 1999 年并于 2005 年 10 月利用火星全球探测者（Mars Global Surveyor）重新分析的假说提出，这些磁带可能表明火星在约 40 亿年前曾具有板块构造活动，当时行星发电机仍在运作，而全球磁场尚未衰亡 \(^{\text{[80]}}\)。
\subsection{地理与地貌}
\begin{figure}[ht]
\centering
\includegraphics[width=8cm]{./figures/6c02165237502845.png}
\caption{标注了各地貌特征的火星地形图，并清晰显示出火星二分性（北部低洼半球与南部高地半球）。} \label{fig_Mars_6}
\end{figure}
尽管更以绘制月球地图而闻名，约翰·海因里希·冯·梅德勒（Johann Heinrich von Mädler）与威廉·贝尔（Wilhelm Beer）是最早的火星地理绘图者。他们首先确认火星的大部分地表特征是永久性的，并更精确地测定了火星的自转周期。1840 年，梅德勒将十年的观测资料整合，绘制出第一幅火星地图 \(^{\text{[81]}}\)。

火星表面的特征名称来源多样。反照率特征以古典神话命名；直径大于约 50 公里的陨石坑以已故科学家、作家及其他对火星研究有贡献的人命名；较小的陨石坑则以全球人口不足 10 万的城镇命名。大型峡谷以不同语言中“火星”或“星”的词汇命名；较小的峡谷则以河流命名 \(^{\text{[82]}}\)。

大型反照率特征保留了许多古老的命名，但常根据对其性质的新认识进行更新。例如，原名为 Nix Olympica（奥林匹斯之雪）的特征现名为 Olympus Mons（奥林匹斯山）\(^{\text{[83]}}\)。从地球上观测，火星表面按反照率差异分为两类区域。较浅色的平原被红铁氧化物丰富的尘埃与沙覆盖，曾被认为是火星的“大陆”，因而被命名为阿拉伯高地（Arabia Terra，阿拉伯之地）或亚马逊平原（Amazonis Planitia，亚马逊平原）。暗色特征则被认为是海洋，因此命名包括 Mare Erythraeum、Mare Sirenum 与 Aurorae Sinus。地球上可见的最大暗区是 Syrtis Major Planum \(^{\text{[84]}}\)。永久的北极冰盖称为 Planum Boreum，南极冰盖称为 Planum Australe \(^{\text{[85]}}\)。

火星赤道由其自转定义，但其本初子午线的位置与地球相同，由人为选定：梅德勒和贝尔在 1830 年制作的第一幅火星地图上选定了一条经线。1972 年“海盗 9 号”（Mariner 9）提供了大量火星影像后，位于子午线海（Sinus Meridiani，“中央海湾”或“子午线海湾”）中的一个小陨石坑（后称 Airy-0）被梅尔顿·E·戴维斯、哈罗德·马苏斯基与热拉尔·德·沃科勒选为 0.0° 经度的标定点，以保持与最初选定的位置一致 \(^{\text{[86][87][88]}}\)。

由于火星没有海洋，因此也不存在“海平面”；必须选取一个零高度面作为参考，即火星的准地球体面（areoid），它与地球的地球体（geoid）相对应 \(^{\text{[89][90]}}\)。零高度被定义为大气压力为 610.5 Pa（6.105 mbar）的高度 \(^{\text{[91]}}\)。该压力对应于水的三相点，大约是地球海平面气压的 0.6\%（0.006 atm）\(^{\text{[92]}}\)。

为制图目的，美国地质调查局将火星表面划分为三十个测图分区，每个分区以其中包含的经典反照率特征命名 \(^{\text{[93]}}\)。2023 年 4 月，《纽约时报》报道了一幅基于“希望号”（Hope）探测器影像更新的火星全球地图 \(^{\text{[94]}}\)。另一幅相关但细节更高的全球火星地图由 NASA 于 2023 年 4 月 16 日发布 \(^{\text{[95]}}\)。
\subsubsection{火山}
\begin{figure}[ht]
\centering
\includegraphics[width=8cm]{./figures/cc4c46ae5220f270.png}
\caption{火星最高的火山——奥林匹斯山（Olympus Mons）的图像。其横跨范围约为 550 公里（340 英里）。} \label{fig_Mars_7}
\end{figure}
广阔的塔尔西斯（Tharsis）高原地区包含若干巨型火山，其中包括盾状火山奥林匹斯山（Olympus Mons）。其火山体宽度超过 600 公里（370 英里）\(^{\text{[96][97]}}\)。由于这座山体规模极其庞大，且其边缘具有复杂结构，因此难以给出确切的高度。从其西北边缘峭壁的山脚至山顶的局部高差超过 21 公里（13 英里）\(^{\text{[97]}}\)，略高于从海底基部测量的冒纳凯阿火山（Mauna Kea）高度的两倍。从亚马孙平原（Amazonis Planitia）出发，向西北方向跨越超过 1,000 公里（620 英里）至火山峰顶的整体海拔变化接近 26 公里（16 英里）\(^{\text{[98]}}\)，约为珠穆朗玛峰（Mount Everest，海拔略高于 8.8 公里〔5.5 英里〕）的三倍。因此，奥林匹斯山是太阳系中最高或第二高的山体；唯一可能更高的是灶神星（Vesta）上的雷亚西尔维亚（Rheasilvia）中央峰，高度约为 20–25 公里（12–16 英里）\(^{\text{[99]}}\)。
\subsubsection{撞击地形}  
火星地形的二元性极为显著：被熔岩流抹平的北部平原，与因古老撞击而密布坑洼的南部高地形成强烈对比。可能在约 40 亿年前，火星北半球曾被一个大小约为地球月球十分之一至三分之二的天体撞击。如果这一假说成立，则火星北半球将是一个尺寸达 10,600 × 8,500 公里（6,600 × 5,300 英里）的撞击盆地，其面积大致相当于欧洲、亚洲与澳大利亚总和，超过乌托邦平原（Utopia Planitia）以及月球南极—艾特肯盆地（South Pole–Aitken Basin），成为太阳系中最大的撞击盆地 \(^{\text{[100][101][102]}}\)。

火星上直径至少 5 公里（3.1 英里）的撞击坑数量达 43,000 个 \(^{\text{[103]}}\)。暴露最明显的大型撞击坑为赫拉斯盆地（Hellas），其直径 2,300 公里（1,400 英里），深达 7,000 米（23,000 英尺），并作为一个明亮的反照率特征从地球上即可清晰观测到 \(^{\text{[104][105]}}\)。其他重要的撞击构造包括阿耳戈里平原（Argyre），其直径约为 1,800 公里（1,100 英里）\(^{\text{[106]}}\)，以及伊西迪斯盆地（Isidis），其直径约为 1,500 公里（930 英里）\(^{\text{[107]}}\)。由于火星的质量和尺寸较小，其受到天体撞击的概率约为地球的一半；但由于火星更靠近小行星带，它受到来自该区域物质撞击的几率更高。火星也更容易遭到短周期彗星的撞击，即那些轨道位于木星轨道以内的彗星 \(^{\text{[108]}}\)。

部分火星撞击坑呈现的形态特征表明地表可能在撞击后变得潮湿 \(^{\text{[109]}}\)。
\subsubsection{构造地点}
\begin{figure}[ht]
\centering
\includegraphics[width=8cm]{./figures/dabff5c6b996abd3.png}
\caption{由“海盗 1 号”（Viking 1）探测器拍摄的水手谷（Valles Marineris）。} \label{fig_Mars_8}
\end{figure}
巨型峡谷水手谷（Valles Marineris，其拉丁语意为“水手号峡谷”，在早期运河地图中亦称为 Agathodaemon \(^{\text{[110]}}\)）全长约 4,000 公里（2,500 英里），深度可达 7 公里（4.3 英里）。其长度相当于欧洲的全长，并横跨火星周长的五分之一。相比之下，地球上的科罗拉多大峡谷（Grand Canyon）仅长 446 公里（277 英里），深近 2 公里（1.2 英里）。水手谷的形成被认为源于塔尔西斯（Tharsis）地区的隆起，使得水手谷区域的地壳发生断陷。2012 年，一项研究提出水手谷不仅仅是一个地堑结构，而可能是一条板块边界，其上发生了约 150 公里（93 英里）的横向位移，这意味着火星可能拥有由两块构成的板块构造体系 \(^{\text{[111][112]}}\)。
\subsubsection{洞穴与孔洞} 
美国宇航局“火星奥德赛号”（Mars Odyssey）轨道器搭载的热发射成像系统（THEMIS）拍摄的图像在阿耳狄西亚山（Arsia Mons）的山坡上揭示了七个可能的洞穴入口 \(^{\text{[113]}}\)。这些洞穴以发现者亲人的名字命名，被统称为“七姐妹”（seven sisters）\(^{\text{[114]}}\)。洞穴入口宽度从 100 到 252 米（328 至 827 英尺）不等，推测深度至少为 73 至 96 米（240 至 315 英尺）。由于大多数洞穴的底部无法被光照及成像到，它们可能远比这些估计值更深，并在地下进一步扩展。“Dena”是唯一例外；其底部可见，并测得深度为 130 米（430 英尺）。这些洞穴内部可能免受微陨石、紫外辐射、太阳耀斑及轰击火星表面的高能粒子的影响，从而形成受保护的环境 \(^{\text{[115][116]}}\)。
\subsubsection{其他特征}
\begin{figure}[ht]
\centering
\includegraphics[width=6cm]{./figures/6aecb1f42844b25f.png}
\caption{艺术家概念图，展示从火星间歇泉喷出的夹沙喷流（由 NASA 发布，艺术家：Ron Miller）。} \label{fig_Mars_9}
\end{figure}
火星间歇泉（Martian geysers，或称二氧化碳喷口 CO\(_2\) jets）被认为是火星南极地区在春季融化期间发生的小规模气体与尘埃喷发地点。“暗沙丘斑点”（dark dune spots）与“蜘蛛”或蜘蛛状地貌（araneiforms）是这类喷发所导致的两种最显著的地表特征 \(^{\text{[117]}}\)。

本节摘自“火星风化层（Martian regolith）”中“§ 大气尘埃（Atmospheric dust）”部分。

轨道上观测到的火星尘暴细节。

火星在没有尘暴时（左图，2001 年 6 月）与发生全球性尘暴时（右图，2001 年 7 月）的比较图像，由火星全球勘测者（Mars Global Surveyor）拍摄。

2018 年 5 月至 9 月期间的火星尘暴光学深度（tau）变化，由火星气候探测仪（Mars Climate Sounder）测量。

尘埃与水云的差异：图像下部中央的黄色云层为大型尘埃云，而其他白色云层为水云。

与地球相比，类似粒径的尘埃会更快从较稀薄的火星大气中沉降。例如，2001 年火星全球性尘暴扬起的尘埃仅在大气中停留约 0.6 年，而皮纳图博火山（Mount Pinatubo）喷发的尘埃在地球大气中约需两年才完全沉降 \(^{\text{[118]}}\)。然而，在当前火星环境条件下，涉及的物质迁移量一般远小于地球。即便是 2001 年火星全球尘暴，其移动的尘埃量折算成均匀沉积层也仅约 3 微米厚，分布在南纬 58°至北纬 58°之间 \(^{\text{[118]}}\)。在两处火星车着陆点的尘埃沉积速度大约为每 100 个火星日（sol）沉积一粒尘埃颗粒的厚度 \(^{\text{[119]}}\)。
\begin{figure}[ht]
\centering
\includegraphics[width=8cm]{./figures/9be5f19566825e95.png}
\caption{暗沙丘斑点} \label{fig_Mars_10}
\end{figure}\begin{figure}[ht]
\centering
\includegraphics[width=8cm]{./figures/9e6fe720e071c2c8.png}
\caption{从轨道上观测到的火星尘暴细节。} \label{fig_Mars_11}
\end{figure}\begin{figure}[ht]
\centering
\includegraphics[width=8cm]{./figures/200ce86596732349.png}
\caption{2001 年 6 月无尘暴的火星（左）与 2001 年 7 月发生全球性尘暴的火星（右），由火星全球勘测者（Mars Global Surveyor）拍摄。} \label{fig_Mars_12}
\end{figure}\begin{figure}[ht]
\centering
\includegraphics[width=8cm]{./figures/35c9fae96a3ee78b.png}
\caption{尘埃云与水云的差异：图像下部中央的黄色云层是大型尘埃云，其他白色云层则是水云。} \label{fig_Mars_13}
\end{figure}
\begin{figure}[ht]
\centering
\includegraphics[width=10cm]{./figures/7e0122ca60745c31.png}
\caption{由“好奇号”（Curiosity）火星车拍摄的火星尘卷风。/一个直径 30 米、高度 800 米的尘卷风。曾观测到高度达 20 公里的尘卷风。/尘卷风在火星表面形成扭曲的深色轨迹。 } \label{fig_Mars_14}
\end{figure}
\subsection{大气}
\begin{figure}[ht]
\centering
\includegraphics[width=8cm]{./figures/f6a34681d67508ea.png}
\caption{} \label{fig_Mars_15}
\end{figure}
火星在约 40 亿年前失去了其磁层 \(^{\text{[120]}}\)，可能是由于多次小行星撞击所致 \(^{\text{[121]}}\)，因此太阳风可直接与火星电离层相互作用，通过剥离外层大气原子而降低大气密度 \(^{\text{[122]}}\)。火星全球勘测者（Mars Global Surveyor）与火星快车号（Mars Express）均探测到火星背后延伸入太空的电离大气粒子尾迹 \(^{\text{[120][123]}}\)，这种大气逃逸现象正由 MAVEN 轨道器进行研究。与地球相比，火星大气极为稀薄。当前火星地表大气压从奥林匹斯山（Olympus Mons）上的 30 Pa（0.0044 psi）低值，到赫拉斯平原（Hellas Planitia）上的超过 1,155 Pa（0.1675 psi）高值不等，平均地表压力约为 600 Pa（0.087 psi）\(^{\text{[124]}}\)。火星上能够达到的最高大气密度，相当于地球表面上方 35 公里（22 英里）处的大气密度 \(^{\text{[125]}}\)。因此火星的平均地表压力仅为地球 101.3 kPa（14.69 psi）的 0.6\%。火星大气的尺度高度约为 10.8 公里（6.7 英里）\(^{\text{[126]}}\)，高于地球的 6 公里（3.7 英里），这是因为火星表面重力仅为地球的约 38\% \(^{\text{[127]}}\)。

火星大气由约 96\% 的二氧化碳、1.93\% 的氩气与 1.89\% 的氮气组成，并含有微量的氧气与水 \(^{\text{[2][128][122]}}\)。大气中含有约 1.5 μm 大小的尘埃颗粒，使得从地表看去火星天空呈黄褐色 \(^{\text{[129]}}\)；由于悬浮的铁氧化物颗粒，它有时也呈粉红色 \(^{\text{[22]}}\)。火星大气中的甲烷浓度从北半球冬季约 0.24 ppb 波动到夏季约 0.65 ppb \(^{\text{[130]}}\)。其寿命估计为 0.6 至 4 年 \(^{\text{[131][132]}}\)，因此其存在意味着必须有活跃的气体来源。甲烷最可能由非生物过程产生，例如含水、二氧化碳与富含橄榄石的岩石之间的蛇纹石化反应，而橄榄石在火星上广泛存在 \(^{\text{[133]}}\)。然而，甲烷也可能由火星生命产生 \(^{\text{[134]}}\)。
\begin{figure}[ht]
\centering
\includegraphics[width=8cm]{./figures/a9dbad45b69fb2ac.png}
\caption{由 MAVEN 在紫外波段观测到的火星逃逸大气（碳、氧与氢）。\(^{\text{[135]}}\)} \label{fig_Mars_16}
\end{figure}
与地球相比，火星大气中较高的二氧化碳浓度与较低的地表气压可能是声音在火星上衰减更为明显的原因；在火星上，自然声源除风以外十分稀少。基于“毅力号”（Perseverance）火星车获取的声学录音，研究者得出结论：当地声速在 240 Hz 以下的频率约为 240 m/s，而在更高频率下约为 250 m/s \(^{\text{[136][137]}}\)。

火星上已观测到极光现象 \(^{\text{[138][139][140]}}\)。由于火星缺乏全球磁场，其极光的类型与分布均不同于地球；与地球极光主要局限于极区不同，火星极光可遍布整颗行星 \(^{\text{[141][142]}}\)。2017 年 9 月，NASA 报道火星表面辐射水平在一次大规模且出乎意料的太阳风暴期间暂时增加至平时的两倍，并伴随出现亮度为此前观测到的 25 倍的极光 \(^{\text{[142][143]}}\)。
\subsubsection{气候}
\begin{figure}[ht]
\centering
\includegraphics[width=8cm]{./figures/8da4b831c9552fc8.png}
\caption{2001 年 7 月无全球尘暴的火星（左）与发生全球尘暴的火星（右），可见不同的水冰云层，由哈勃太空望远镜（Hubble Space Telescope）拍摄。} \label{fig_Mars_17}
\end{figure}
\begin{figure}[ht]
\centering
\includegraphics[width=8cm]{./figures/8546f9a4ee0dacc6.png}
\caption{火星北极（左）与南极（右）在北半球与南半球夏季之间二氧化碳冰（非水冰）覆盖变化的渲染图。} \label{fig_Mars_18}
\end{figure}
火星具有类似地球的季节变化，在其北半球与南半球之间交替出现。此外，相较于地球轨道，火星的轨道离心率更大：当南半球为夏季、北半球为冬季时，火星位于近日点；当南半球为冬季、北半球为夏季时，火星位于远日点。其结果是，南半球的季节变化更为极端，而北半球的季节则比本应出现的情况更为温和。南半球的夏季温度可比北半球对应季节高出多达 30 °C（54 °F）\(^{\text{[144]}}\)。

火星地表温度在赤道夏季的最高可达约 35 °C（95 °F），最低可达约 −110 °C（−166 °F）\(^{\text{[16]}}\)。这一极大的温差源于火星稀薄的大气无法储存多少太阳热量、大气压仅为地球大气压的约 1\%，以及火星土壤的低热惯量 \(^{\text{[145]}}\)。火星距离太阳的平均距离是地球的 1.52 倍，因此其接收到的阳光量仅为地球的 43\% \(^{\text{[146][147]}}\)。

火星拥有太阳系中最大的沙尘暴，其风速可超过 160 km/h（100 mph）。其规模可从局部区域的小型风暴到覆盖整颗行星的全球性沙尘暴。它们通常在火星最接近太阳时出现，并已被证明会提升全球气温 \(^{\text{[148]}}\)。季节变化还会使极冠形成干冰覆盖层 \(^{\text{[149]}}\)。
\subsection{水文}
\begin{figure}[ht]
\centering
\includegraphics[width=8cm]{./figures/8605fe891fd8f42e.png}
\caption{火星含有水，但主要以被尘埃覆盖的极地冰层形式存在，如本图所示。} \label{fig_Mars_19}
\end{figure}
尽管火星含有较大量的水，但其中大部分以被尘埃覆盖的水冰形式存在于火星极地冰冠之中 \(^{\text{[150][151][152][153][85]}}\)。如果南极冰冠中的水冰全部融化，其体积足以在火星大部分表面形成一层约 11 米（36 英尺）深的水层 \(^{\text{[154]}}\)。

由于火星的大气压不足地球的 1\%，液态水无法在其地表长期存在 \(^{\text{[155]}}\)。只有在最低洼的区域，气压与温度才可能在短时间内满足液态水存在的条件 \(^{\text{[47][156]}}\)。

虽然火星大气中的水含量很少，但仍足以形成水冰云，以及不同形式的雪与霜，这些现象常与二氧化碳干冰雪混合出现 \(^{\text{[157]}}\)。
\subsubsection{过去的水圈}
\begin{figure}[ht]
\centering
\includegraphics[width=8cm]{./figures/cc4f1bb7411d46a7.png}
\caption{艺术家笔下四十亿年前的火星想象图。} \label{fig_Mars_20}
\end{figure}
火星上可见的各种地貌强烈表明液态水曾存在于其地表。巨大的线性冲蚀带，称为溢流槽（outflow channels），在约 25 处横贯火星表面。一般认为，这些地貌记录了地下含水层中大量水突发释放所造成的侵蚀，但也有部分结构被假设可能由冰川或熔岩作用形成 \(^{\text{[158][159]}}\)。其中一处规模较大的例子——马阿迪姆谷（Ma'adim Vallis）长达 700 公里（430 英里），远大于地球上的科罗拉多大峡谷，其宽度达 20 公里（12 英里），部分区域深达 2 公里（1.2 英里）。它被认为是在火星早期由流动水切割而成 \(^{\text{[160]}}\)。这些溢流槽中最年轻的一部分被认为仅在几百万年前形成 \(^{\text{[161]}}\)。

在火星的其他地区，尤其是火星表面最古老的区域，遍布更为细密、树状分支状的谷网。这些谷地形态及其分布强烈暗示它们是由降水产生的径流在火星早期历史中切割而成。地下水流动与地下水侧蚀（groundwater sapping）可能在部分谷网中发挥次要作用，但几乎所有情况下，最根本的侵蚀来源应为降水 \(^{\text{[162]}}\)。

在陨石坑和峡谷壁上，有数千处与地球类似的沟槽（gullies）。这些沟槽多见于南半球高地，其方向朝向赤道，并全部位于纬度 30° 以极向一侧。许多研究者认为其形成过程涉及液态水，可能源自冰融化 \(^{\text{[163][164]}}\)，而另一些人则提出可能由二氧化碳霜或干尘移动形成 \(^{\text{[165][166]}}\)。尚未发现因风化而部分退化的沟槽，也未发现覆盖于其上的撞击坑，说明这些地貌非常年轻，并可能仍处于活动状态 \(^{\text{[164]}}\)。其他地质特征，如保存在陨石坑中的三角洲与冲积扇，则进一步支持火星早期某一或数个阶段出现过更温暖、更潮湿的环境 \(^{\text{[167]}}\)。此类环境必然要求火星表面大面积存在陨石坑湖泊，而矿物学、沉积学与地貌学的独立证据均支持这一点 \(^{\text{[168]}}\)。更多证据来自于特定矿物（如赤铁矿和针铁矿）的发现，这些矿物往往在水的参与下形成 \(^{\text{[169]}}\)。
\subsubsection{水的观测历史与相关证据的发现}
\begin{figure}[ht]
\centering
\includegraphics[width=8cm]{./figures/2007889d13f8e0f2.png}
\caption{如本图中这座名为柯罗廖夫（Korolev）的极地陨坑所示，火星表层的水冰在某些区域清晰可见，此为其三维投影图像。} \label{fig_Mars_21}
\end{figure}
2004 年，“机遇号”（Opportunity）探测器检测到黄钾铁矾（jarosite）这一矿物。它仅能在酸性水的存在下形成，因此表明火星上曾经存在过水 \(^{\text{[170][171]}}\)。2007 年，“勇气号”（Spirit）探测器发现高浓度二氧化硅沉积，显示火星过去具有潮湿环境；2011 年 12 月，“机遇号”再次在火星表面发现了石膏，这同样是一种在有水条件下形成的矿物 \(^{\text{[172][173][174]}}\)。估计火星上地幔上部所含的水量——以火星矿物中所含氢氧根离子的形式表示——在水含量为 50–300 ppm 的范围内，与地球相当甚至更高，此量足以在火星表面形成 200–1,000 米（660–3,280 英尺）深的全球性水层 \(^{\text{[175][176]}}\)。

2013 年 3 月 18 日，NASA 报告称“好奇号”（Curiosity）火星车的仪器在多份岩石样品中发现矿物水合现象，推测为水合硫酸钙。这些样品包括破裂的“Tintina”岩与“Sutton Inlier”岩碎片，以及诸如“Knorr”岩与“Wernicke”岩等其他岩石中的脉体与结核 \(^{\text{[177][178]}}\)。使用火星车 DAN 仪器的分析表明，地表以下 60 厘米（24 英寸）范围内存在地下水，其含水量高达 4\%；这些测量是在火星车从布拉德伯里着陆点（Bradbury Landing）驶向格伦尔格（Glenelg）地形中黄刀湾（Yellowknife Bay）区域的过程中获得的 \(^{\text{[177]}}\)。2015 年 9 月，NASA 宣布基于光谱仪对坡面暗色区域的分析，发现了关于重复坡线（recurring slope lineae）中水合盐卤（hydrated brine）流动的强有力证据 \(^{\text{[179][180][181]}}\)。这些条纹在火星夏季向坡下流动，此时温度高于 −23 °C；在更低温度下则冻结 \(^{\text{[182]}}\)。这些观测支持了早先基于其形成时间与增长速率提出的假说，即这些暗条纹源于地表下方水的流动 \(^{\text{[183]}}\)。然而后续研究提出，这些坡线可能是干燥的颗粒流动，其形成过程中水的作用至多有限 \(^{\text{[184]}}\)。目前关于火星地表液态水的存在、范围与作用仍缺乏定论 \(^{\text{[185][186]}}\)。

研究者推测，火星北部的低洼平原曾被一片深达数百米的海洋所覆盖，尽管这一理论仍存在争议 \(^{\text{[187]}}\)。2015 年 3 月，科学家提出这一古海洋可能与地球的北冰洋规模相当。这一结论基于当代火星大气中氕与氘的比值与地球的对应比值的比较。火星大气中的氘含量（D/H = 9.3 ± 1.7 × 10\(^{-4}\)）是地球（D/H = 1.56 × 10\(^{-4}\)）的五至七倍，这表明古代火星曾具有显著更高的水量。此前，“好奇号”（Curiosity）曾在盖尔陨坑（Gale Crater）检测到较高的氘比值，尽管不足以明确证明海洋曾经存在。然而，有科学家警告称这些结果尚未得到进一步确认，并指出当前的火星气候模型尚未显示火星过去具备足够温暖的条件以支持液态水体的长期存在 \(^{\text{[188]}}\)。在北极冰冠附近，火星快车号（Mars Express）发现宽 81.4 公里（50.6 英里）的柯罗廖夫陨坑（Korolev Crater）中约储存有 2,200 立方公里（530 立方英里）的水冰 \(^{\text{[189]}}\)。

2016 年 11 月，NASA 报告在乌托邦平原（Utopia Planitia）地区发现大量地下冰，其水量估计相当于苏必利尔湖（Lake Superior）的水量（12,100 立方公里 \(^{\text{[190]}}\)）\(^{\text{[191][192]}}\)。在 2018 至 2021 年期间，“ExoMars 痕量气体轨道器”（Trace Gas Orbiter）在水手谷（Valles Marineris）峡谷系统中探测到疑似水的迹象，可能为地下冰的存在 \(^{\text{[193]}}\)。
\subsection{轨道运动}
\begin{figure}[ht]
\centering
\includegraphics[width=8cm]{./figures/620ad9bc1db5b447.png}
\caption{火星及太阳系内侧行星的轨道。} \label{fig_Mars_22}
\end{figure}
火星与太阳的平均距离约为 2.3 亿公里（1.43 亿英里），其公转周期为 687 个地球日。火星的太阳日（sol）仅比地球日稍长，为 24 小时 39 分 35.244 秒 \(^{\text{[194]}}\)。一个火星年等于 1.8809 个地球年，即 1 年 320 天 18.2 小时 \(^{\text{[2]}}\)。从地球转移到火星所需跨越的引力势差及相应的变轨速度（delta-v）是所有行星间转移中仅次于地球—金星转移的第二低 \(^{\text{[195][196]}}\)。

火星自转轴相对其轨道平面的倾角为 25.19°，与地球的自转轴倾角十分相近 \(^{\text{[2]}}\)。因此火星亦存在类似地球的季节变化，只是由于火星轨道周期更长，其季节持续时间几乎是地球的两倍。在当今时代，火星北极的指向大致接近天津四（Deneb）\(^{\text{[21]}}\)。

火星的轨道偏心率约为 0.09，相对明显；在太阳系的其他七颗行星中，只有水星的轨道偏心率更大。已知火星过去的轨道要圆得多。例如约 135 万年前，火星的轨道偏心率仅约 0.002，远小于今日地球的偏心率 \(^{\text{[197]}}\)。火星轨道偏心率的变化周期约为 96,000 个地球年，相比之下地球为 100,000 年 \(^{\text{[198]}}\)。

火星与地球的最近接（冲日）发生在 779.94 天的会合周期中。需要注意，这与“合日”（conjunction）不同，后者是地球与火星分处太阳两侧呈一直线。火星两次冲日间的时间即为其会合周期，平均为 780 天；但实际间隔可能在 764 至 812 天之间变化 \(^{\text{[198]}}\)。由于两者轨道均为椭圆形，最近距离可在约 5,400 万至 1.03 亿公里（3,400 万至 6,400 万英里）之间变化，对应角直径也随之变化 \(^{\text{[199]}}\)。在最远时，火星与地球的距离可达 4.01 亿公里（2.49 亿英里）\(^{\text{[200]}}\)。火星每约 2.1 年从地球看会出现一次冲日。在 2003、2018 与 2035 年，火星在近日点附近发生冲日，而 2020 与 2033 年的冲日事件则特别接近于典型的近近日点冲日 \(^{\text{[201][202][203]}}\)。
\begin{figure}[ht]
\centering
\includegraphics[width=6cm]{./figures/2baf2d2fd195eea0.png}
\caption{2020 年火星冲日期间，通过一台 16 英寸业余望远镜观测到的火星。} \label{fig_Mars_23}
\end{figure}
火星的平均视星等为 +0.71，其标准差为 1.05 \(^{\text{[19]}}\)。由于火星轨道具有偏心率，其在与太阳相对冲期间的视星等可在约 −3.0 至 −1.4 之间变化 \(^{\text{[204]}}\)。当火星接近远日点并与太阳处于合日位置时，其最暗，视星等为 +1.86 \(^{\text{[19]}}\)。在最亮时，火星（与木星一起）在视亮度上仅次于金星 \(^{\text{[19]}}\)。火星通常呈显著的黄色、橙色或红色。当其最远离地球时，距离超过其最近时的七倍。火星通常每隔约 15 年或 17 年会有一到两次距离地球足够近，以达到特别好的观测条件 \(^{\text{[205]}}\)。地基光学望远镜的分辨率受地球大气限制，即便在火星与地球最近时，通常也只能解析出约 300 公里（190 英里）尺度的地表特征 \(^{\text{[206]}}\)。

当火星接近冲日时，会出现逆行运动，即相对于背景恒星呈现向后的弧形运动。这段逆行运动持续约 72 天，而火星的视亮度峰值出现在这一区间的中段 \(^{\text{[207]}}\)。
\subsection{卫星}
\begin{figure}[ht]
\centering
\includegraphics[width=6cm]{./figures/572b91d0c563f574.png}
\caption{从火星上看到的火卫二（Deimos）和火卫一（Phobos），其视大小与从地球看到的、更大且实际尺寸也更大的月球作比较。如果这两颗卫星距离火星与月球距离地球相同，它们在火星天空中将只会呈现为微弱的星状亮点。} \label{fig_Mars_24}
\end{figure}
\begin{figure}[ht]
\centering
\includegraphics[width=6cm]{./figures/93860f6dcfe55ddd.png}
\caption{按比例绘制的火星天然与人造卫星轨道图，其中最远的火卫二（Deimos）轨道半径为 23,460 公里（14,580 英里）。} \label{fig_Mars_25}
\end{figure}
火星拥有两颗相对较小（相较于地球）的天然卫星：火卫一（Phobos，直径约 22 公里〔14 英里〕）与火卫二（Deimos，直径约 12 公里〔7.5 英里〕），它们分别在距离火星 9,376 公里（5,826 英里）与 23,460 公里（14,580 英里）的轨道上运行。两颗卫星的起源尚不明确，但一种广为接受的理论认为它们是被火星捕获的小行星 \(^{\text{[208]}}\)。

这两颗卫星均由阿萨夫·霍尔（Asaph Hall）于 1877 年发现，并以希腊神话人物命名：Phobos（恐慌与惊惧之神）与 Deimos（恐怖与害怕之神），他们是战神阿瑞斯（Ares）的双胞胎随从，陪伴其出征 \(^{\text{[209]}}\)。火星在罗马神话中对应阿瑞斯。在现代希腊语中，这颗行星仍保留其古名 Άρης（Aris）\(^{\text{[101]}}\)。

从火星表面看，火卫一与火卫二的运行方式与地球的月球显著不同。火卫一从西方升起，东方落下，并在短短 11 小时后再次升起。火卫二因其轨道距离仅略高于同步轨道（同步轨道上卫星的公转周期与行星自转周期一致），因此如预期般从东方升起，但速度缓慢。由于火卫一轨道低于同步高度，火星的潮汐力正逐渐降低其轨道高度。在约 5,000 万年后，它可能会撞击火星表面或解体并形成行星环 \(^{\text{[210]}}\)。

两颗卫星的起源仍未有定论。其低反照率与碳质球粒陨石样的成分常被视为与小行星相似，从而支持“捕获假说”。火卫一不稳定的轨道似乎也暗示其可能是近期被捕获的。然而这两颗卫星均具有接近赤道的圆形轨道，而捕获天体通常不会呈现此特征，且捕获所需的动力学条件也极为复杂。另一种可能性是它们在火星早期历史中就地聚积而成，但若其成分确实类似小行星而非火星本身，则该理论难以成立 \(^{\text{[211]}}\)。火星可能还存在尚未发现的小型卫星，其直径小于 50 至 100 米（160 至 330 英尺），且预计在火卫一与火卫二之间存在一条尘埃环 \(^{\text{[212]}}\)。

关于两颗卫星起源的第三种可能性是涉及第三天体或撞击破碎过程。近期证据表明火卫一内部高度多孔 \(^{\text{[213]}}\)，其成分可能主要由层状硅酸盐（phyllosilicates）及其他已知存在于火星的矿物组成 \(^{\text{[214]}}\)。这些证据指向火卫一可能源自火星表面一次大型撞击所喷射出的物质，再在火星轨道上重新聚积而成，其形成机制类似主流的地球月球起源理论。尽管两颗卫星在可见光与近红外（VNIR）光谱上与外侧小行星带天体相似，但火卫一的热红外光谱据报道与任何类型的球粒陨石均不一致 \(^{\text{[214]}}\)。另一种可能性是火卫一与火卫二本是更早期的一颗古老卫星的碎片，该卫星由火星一次大型撞击产生，再由于一次较新的撞击而被破坏 \(^{\text{[215]}}\)。
\subsection{人类的观测与探索}
火星的观测史以火星冲日为标志，因为在冲日期间火星最接近地球，因此最易观测；这一现象每隔数年发生一次。更值得注意的是火星近日点冲日，此时火星接近其近日点，使其与地球的距离更近 \(^{\text{[201]}}\)。
\subsubsection{古代观测}  
古代苏美尔人将火星称为尼格尔（Nergal），即战争与瘟疫之神。在苏美尔时期，尼格尔只是一个影响力较小的次级神祇，但在后期，他的主要崇拜中心位于尼尼微（Nineveh）城 \(^{\text{[216]}}\)。在美索不达米亚文献中，火星被称为“审判亡者命运之星” \(^{\text{[217]}}\)。古埃及天文学家也记录了火星作为夜空游移天体的存在，并在公元前 1534 年左右已经熟悉其逆行运动 \(^{\text{[218]}}\)。至新巴比伦帝国时期，巴比伦天文学家定期记录行星位置并系统性观测它们的行为。他们知道火星在 79 年中经历 37 次会合周期，或穿越黄道带 42 圈。他们还发明了用于对预测行星位置进行小幅修正的算术方法 \(^{\text{[219][220]}}\)。在古希腊，火星被称为 Πυρόεις \(^{\text{[221]}}\)，常见称呼则为阿瑞斯（Ares）。罗马人将其命名为 Mars，以他们的战争之神为名，象征其剑与盾的图腾 \(^{\text{[222]}}\)。

公元前四世纪，亚里士多德注意到火星在一次掩星过程中被月球遮掩，这表明火星距离地球比月球更远 \(^{\text{[223]}}\)。居住在亚历山大的希腊天文学家托勒密 \(^{\text{[224]}}\) 试图解释火星轨道运动的问题。他的模型与整体天文学著作被收录于后来称作《天文学大成》（Almagest，源自阿拉伯语“最伟大之书”）的多卷本著作中，并在随后的十四个世纪成为西方天文学的权威著作 \(^{\text{[225]}}\)。古代中国的文献亦证实，早在公元前四世纪，火星已为中国天文学家所知 \(^{\text{[226]}}\)。在东亚文化中，火星传统上依据五行体系被称为“火星” \(^{\text{[227][228][229]}}\)。
\subsubsection{近代早期观测}
\begin{figure}[ht]
\centering
\includegraphics[width=6cm]{./figures/e8b0b841cb25c834.png}
\caption{约翰内斯·开普勒（Johannes Kepler）于《新天文学》（Astronomia Nova，1609）中绘制的火星在若干次视逆行期间的地心轨迹示意图。} \label{fig_Mars_26}
\end{figure}
1609 年，约翰内斯·开普勒（Johannes Kepler）发表了对火星轨道长达十年的研究 \(^{\text{[230]}}\)，并利用第谷·布拉赫（Tycho Brahe）测得的火星日视差对火星的相对距离进行了初步计算 \(^{\text{[231]}}\)。基于布拉赫对火星的观测数据，开普勒推断火星绕太阳运行的轨道并非圆形，而是椭圆。此外，开普勒还证明火星在接近太阳时加速，在远离太阳时减速，这一现象后来被物理学家解释为角动量守恒的结果 \(^{\text{[232:433–437]}}\)。
\begin{figure}[ht]
\centering
\includegraphics[width=6cm]{./figures/692b099095f6b294.png}
\caption{1659 年 11 月 28 日由克里斯蒂安·惠更斯（Christiaan Huygens）观测并绘制的火星地貌最早记录（手稿 K）\(^{\text{[233]}}\)。图中所示特征很可能为大瑟提斯平原（Syrtis Major Planum）。惠更斯根据这些地貌的移动计算出了火星的自转周期。} \label{fig_Mars_27}
\end{figure}
1610 年，意大利天文学家伽利略·伽利莱（Galileo Galilei）首次使用望远镜进行天文观测，其中包括对火星的观测 \(^{\text{[234]}}\)。利用望远镜，人们再次测量火星的日视差，以期确定日地距离。乔瓦尼·多梅尼科·卡西尼（Giovanni Domenico Cassini）于 1672 年首次成功实施了这一测量。然而，由于仪器质量有限，早期的视差测量受到严重限制 \(^{\text{[235]}}\)。唯一一次被观察到的金星掩火星事件发生于 1590 年 10 月 13 日，由迈克尔·梅斯特林（Michael Maestlin）在海德堡观测到 \(^{\text{[236]}}\)。
\subsubsection{火星“运河”}
\begin{figure}[ht]
\centering
\includegraphics[width=8cm]{./figures/ae91ead54de6e9ce.png}
\caption{1962 年由美国航空航图与信息中心（U.S. Aeronautical Chart and Information Center）发布的火星地图，其中显示了蜿蜒穿过火星地表的“运河”。当时，运河的存在仍高度具争议性，因为在此之前尚未获得任何火星的近距离图像（直到 1965 年水手 4 号的飞掠）。} \label{fig_Mars_28}
\end{figure}
到了 19 世纪，望远镜的分辨率已经足以辨认火星表面的特征。1877 年 9 月 5 日，火星发生了一次近日点冲日。意大利天文学家乔瓦尼·夏帕雷利（Giovanni Schiaparelli）在米兰使用一台 22 厘米（8.7 英寸）望远镜绘制了第一幅火星的详细地图。这些地图显著地标注了他称为 canali 的特征，而这些特征（除自然峡谷水手谷 Valles Marineris 之外的所有部分）后来被证明是视觉错觉。他认为这些 canali 是火星表面又长又直的线条，并以地球上著名河流为其命名。他所用的词 canali 意为“水道”或“槽沟”，但在英语中却被大众误译为 “canals”（人工运河）\(^{\text{[237][238]}}\)。

受到这些观测的影响，东方学家珀西瓦尔·洛厄尔（Percival Lowell）建立了一座天文台，配备了 30 与 45 厘米（12 与 18 英寸）望远镜。这座天文台用于 1894 年最后一次良好的火星观测机会，以及随后的数次较不理想的冲日观测。他出版了多部关于火星及其生命可能性的著作，对公众产生了巨大影响 \(^{\text{[239][240]}}\)。这些 canali 也被其他天文学家独立观测到，如尼斯的亨利·约瑟夫·佩罗坦（Henri Joseph Perrotin）与路易·托隆（Louis Thollon），他们使用了当时规模最大的望远镜之一 \(^{\text{[241][242]}}\)。

火星季节性变化（包括极冠缩小以及火星夏季出现的暗色区域）与这些所谓的运河共同促成了关于火星生命的种种猜测，人们长期相信火星上存在广阔的海洋与植被。但随着更大口径望远镜的使用，被观测到的直线状 canali 越来越少。1909 年，安东尼亚迪（Antoniadi）使用一台 84 厘米（33 英寸）望远镜观测火星时，只看到不规则的地表图案，并未发现任何 canali \(^{\text{[243]}}\)。
\subsubsection{首次探索}
\begin{figure}[ht]
\centering
\includegraphics[width=6cm]{./figures/7e640647e3592c77.png}
\caption{水手 4 号（Mariner 4）拍摄并传回的首两张来自另一颗行星（火星）的图像 \(^{\text{[244]}}\)。水手 4 号是第二艘抵达火星的航天器，也是第一艘能够传输数据的航天器（1965 年 7 月 15 日）。其传回的数据否定了火星宜居及存在“运河”的假设。} \label{fig_Mars_29}
\end{figure}
地球发往火星的第一艘航天器是苏联的“火星 1 号”（Mars 1），其于 1963 年飞掠火星，但在途中失去通讯。随后是 NASA 的水手 4 号（Mariner 4），它成为第一艘成功从火星传回数据的航天器；该探测器于 1964 年 11 月 28 日发射，并于 1965 年 7 月 15 日最接近火星。水手 4 号探测到微弱的火星辐射带，其强度约为地球的 0.1\%，并拍摄了深空中首次来自另一颗行星的图像 \(^{\text{[245]}}\)。

随着航天器在 1960 年代与 1970 年代访问火星，人们此前对火星的诸多想象被彻底推翻。在“海盗号”（Viking）生命探测实验结果公布后，“死寂星球”这一假说被广泛接受 \(^{\text{[246]}}\)。水手 9 号（Mariner 9）与海盗号任务提供的数据使得更加精确的火星地图得以绘制。
\subsubsection{探索的再度复兴} 
在 1982 年海盗 1 号停止工作之后，直到 1997 年之间，火星仅被三次探测但均未成功：其中两次探测器飞掠火星但未建立通讯（火卫一号 Phobos 1，1988；火星观测者号 Mars Observer，1993），另一次（火卫二号 Phobos 2，1989）在进入轨道后到达目标火卫二前发生故障。

1997 年，“火星探路者号”（Mars Pathfinder）成为继月球之后首个成功的火星漫游者任务，并与“火星全球探勘者号”（Mars Global Surveyor，运行至 2006 年末）一道，开启了持续至今的火星机器人探测时代。该任务生成了完备、极高分辨率的火星地形、磁场与地表矿物分布图 \(^{\text{[247]}}\)。

从这些任务开始，一系列改进型无人航天器——包括轨道器、着陆器与漫游者——陆续被送往火星。NASA（美国）、JAXA（日本）、ESA、英国、ISRO（印度）、俄罗斯国家航天集团（Roscosmos）、阿拉伯联合酋长国，以及中国国家航天局（CNSA）均成功执行了相关探测任务，以研究火星的地表、气候与地质 \(^{\text{[248]}}\)，从而揭示火星水圈的演化历史、其动力过程，以及古代生命可能留下的痕迹。
\begin{figure}[ht]
\centering
\includegraphics[width=8cm]{./figures/90458c73a47aff61.png}
\caption{“毅力号”（Perseverance）火星车与“机智号”（Ingenuity）直升机在莱特兄弟飞行场（Wright Brothers Field）拍摄的合照，2021 年。} \label{fig_Mars_30}
\end{figure}
\subsubsection{当前任务}
截至 2023 年，火星上共有十个正常运行的航天器。

其中八个在轨运行：2001 火星奥德赛号（2001 Mars Odyssey）、火星快车号（Mars Express）、火星勘测轨道器（Mars Reconnaissance Orbiter）、MAVEN、ExoMars 痕量气体轨道器（Trace Gas Orbiter）、“希望号”（Hope orbiter）以及“天问一号”轨道器 \(^{\text{[249][250]}}\)。

另有两个位于火星表面：“火星科学实验室”——“好奇号”（Curiosity）漫游者，以及“毅力号”（Perseverance）漫游者 \(^{\text{[251]}}\)。

收集整理的火星地图可在包括 Google Mars 在内的网站在线浏览。NASA 还提供两个在线工具：“Mars Trek”，可利用 50 年来火星探测的数据对火星进行可视化；以及“Experience Curiosity”，能够以 3D 形式模拟驾驶“好奇号”在火星上行进 \(^{\text{[252][253]}}\)。
\subsubsection{未来任务}  
计划中的火星任务包括：

\begin{itemize}
\item NASA 的 EscaPADE 航天器，于 2025 年 11 月 13 日发射 \(^{\text{[254][255][256]}}\)；  
\item 旨在寻找古代生命证据的 Rosalind Franklin 漫游者任务，原定于 2018 年发射，但多次延期，目前最早的发射日期推迟至 2028 年 \(^{\text{[257][258][259]}}\)，并于 2024 年在追加经费支持下重新启动 \(^{\text{[260]}}\)；  
\item NASA 与欧洲航天局（ESA）联合的火星样本返回任务的当前构想，计划于 2026 年发射 \(^{\text{[261][262]}}\)；  
\item 中国的“天问三号”（Tianwen-3）样本返回任务，预计于 2028 年或 2030 年发射 \(^{\text{[263]}}\)。
\end{itemize}
截至 2024 年 2 月，这类火星任务遗留的碎片总量已超过七吨，其中大部分由坠毁或失效的航天器及其抛弃的部件构成 \(^{\text{[264][265]}}\)。

2024 年 4 月，NASA 选定多家公司开展研究，以提供商业服务，进一步促进火星机器人科学探索的发展。重点领域包括建立通信系统、运载有效载荷以及提供地表成像服务 \(^{\text{[266]}}\)。
\subsection{宜居性与居住}
\begin{figure}[ht]
\centering
\includegraphics[width=8cm]{./figures/925397080cdbe775.png}
\caption{在切亚瓦瀑布（Cheyava Falls）古老的泥质化石沉积物图像中，可见层理、结核、反应前缘以及有机物探测信号，这些发现被认为是迄今最有希望揭示古代火星生命遗迹的证据。} \label{fig_Mars_31}
\end{figure}
19 世纪晚期，天文学界普遍认为火星具有支持生命的条件，包括存在氧气与水 \(^{\text{[267]}}\)。然而，1894 年，利克天文台（Lick Observatory）的 W. W. Campbell 观测火星后发现，“如果火星大气中存在水汽或氧，其含量也小到当时的光谱仪无法探测” \(^{\text{[267]}}\)。该结果与当时许多测量相矛盾，因此未被广泛接受。1909 年，Campbell 与 V. M. Slipher 使用性能更佳的仪器重复了研究，仍得到相同的结果。直到 1925 年 W. S. Adams 的工作最终确认这一结论，关于火星具备类地宜居性的神话才被彻底打破 \(^{\text{[267]}}\)。然而，即使在 1960 年代，仍有文章发表讨论火星生物学，并将火星季节变化归因于生命活动，而非其他物理过程 \(^{\text{[268]}}\)。

当前对行星宜居性的理解——即一颗行星形成适合生命起源的环境条件的能力——强调行星表面需存在液态水。通常这要求行星轨道位于宜居带内，对于太阳而言，宜居带被认为从地球轨道内部延伸至火星轨道附近 \(^{\text{[269]}}\)。在近日点时，火星短暂进入宜居带，但其稀薄的大气（低压）使液态水无法在大范围、长时间存在。火星过去存在液态水流动的证据表明其曾具备一定的宜居潜力。但近来的证据显示，火星地表若存在水，可能咸度和酸度都过高，不足以支持典型陆地生命 \(^{\text{[270]}}\)。

火星的环境条件对有机生命的维持构成巨大挑战：其表面热量传输能力极弱；由于缺乏全球磁层，难以抵御太阳风轰击；其大气压不足以让水以液态长期存在（水会直接升华为气态）。火星几乎——或可能完全——已经停止地质活动；火山活动的终止似乎阻断了地表与行星内部之间化学与矿物质的循环 \(^{\text{[271]}}\)。

证据表明，火星曾经比如今显著更具宜居性，但是否曾存在生命仍然未知。20 世纪 70 年代中期的“海盗号”（Viking）探测器在着陆地点进行了旨在探测火星土壤中微生物的实验，并获得了阳性结果，包括在加入水和营养物后二氧化碳产量的暂时上升。这一生命迹象后来被科学家提出质疑，因此仍存在持续争论；NASA 科学家 Gilbert Levin 坚持认为“海盗号”可能已经发现了生命 \(^{\text{[272]}}\)。2014 年对火星陨石 EETA79001 的分析发现其中含有氯酸盐、高氯酸盐和硝酸盐离子，其含量之高表明这些离子在火星上可能广泛存在。紫外线和 X 射线辐射会将氯酸盐和高氯酸盐转化为更具反应性的含氧氯化物，因此任何有机分子若想幸存，必须埋藏于地表以下 \(^{\text{[273]}}\)。

火星轨道器检测到少量甲烷与甲醛，这两种化合物在火星大气中会迅速分解，因此被一些研究者认为可能是生命存在的证据 \(^{\text{[274][275]}}\)。但另一种解释是，这些化合物可能由火山或其他地质过程（如蛇纹石化作用）不断补充 \(^{\text{[133]}}\)。由陨石撞击生成的冲击玻璃也已在火星陨坑表面被发现，在地球上，这类玻璃能够保存生命迹象 \(^{\text{[276][277]}}\)。因此，如果火星曾在相关地点存在生命，那么火星陨坑中的冲击玻璃也可能保存其痕迹 \(^{\text{[278][279][280]}}\)。

2024 年 6 月在火星上发现的切亚瓦瀑布（Cheyava Falls）岩石被 NASA 指定为“潜在生物特征”（potential biosignature），并已由“毅力号”（Perseverance）火星车钻取岩芯，以用于未来取回地球进行进一步分析。尽管这一发现极具吸引力，但依据当前数据尚无法对其是否具有生物或非生物起源作出最终结论 \(^{\text{[281]}}\)。
\subsubsection{载人任务提案}
\begin{figure}[ht]
\centering
\includegraphics[width=8cm]{./figures/339681bf3d4dc3a9.png}
\caption{NASA 的一种原位资源利用（ISRU）系统构想：由自主机器人开挖并处理火星土壤，以提取探索任务所需的水资源。} \label{fig_Mars_32}
\end{figure}
多个火星载人任务计划曾被提出，但都尚未实现。《2017 年 NASA 授权法案》要求 NASA 研究 2030 年代早期执行载人火星任务的可行性；最终报告得出结论认为此举不可行 \(^{\text{[282][283]}}\)。此外，2021 年，中国也宣布计划于 2033 年执行载人火星任务 \(^{\text{[284]}}\)。SpaceX 等私人企业亦提出了将人类送往火星并最终在火星上定居的计划 \(^{\text{[285]}}\)。截至 2024 年，SpaceX 正在推进“星舰”（Starship）运载系统的研发，其目标即为实现火星殖民。依据 2024 年 4 月与公司内部分享的计划，埃隆·马斯克（Elon Musk）设想在未来二十年内开始建立火星殖民地。这将通过大规模制造星舰、并在初期由地球补给与火星原位资源利用（ISRU）共同支持，直至殖民地实现完全自给自足 \(^{\text{[286]}}\)。任何未来的火星载人任务都可能选择每隔 26 个月出现一次的火星最佳发射窗口。火卫一（Phobos）也被提议作为建设太空电梯的锚点 \(^{\text{[287]}}\)。除各国航天机构与商业航天公司外，火星学会（Mars Society）\(^{\text{[288]}}\) 与行星学会（The Planetary Society）\(^{\text{[289]}}\) 等民间组织亦积极倡导开展火星载人任务。
\subsection{在文化中}
\begin{figure}[ht]
\centering
\includegraphics[width=6cm]{./figures/39009749c415b955.png}
\caption{H. G. 威尔斯（H. G. Wells）于 1897 年出版的《世界大战》（The War of the Worlds）描绘了虚构的火星人入侵地球的情节。} \label{fig_Mars_33}
\end{figure}
火星以罗马战神命名（希腊语为阿瑞斯 Ares），但古希腊天文学家也曾将其与半神赫拉克勒斯（罗马的赫拉克勒斯 Hercules）联系起来，亚里士多德对此有详尽记载 \(^{\text{[290]}}\)。火星与战争的关联至少可追溯至巴比伦天文学，在那里火星被命名为尼格尔（Nergal），即战争与毁灭之神 \(^{\text{[291][292]}}\)。这种象征一直延续至近代，例如霍尔斯特（Gustav Holst）的管弦乐组曲《行星组曲》（The Planets），其著名的第一乐章便将火星称作“战争使者”（the bringer of war）\(^{\text{[293]}}\)。火星的符号为右上方带长矛的圆形，该符号亦被用于代表男性性别 \(^{\text{[294]}}\)。这一符号可追溯至至少 11 世纪，而其潜在的前身符号在希腊的《牛首纸草文稿》（Oxyrhynchus Papyri）中也曾出现 \(^{\text{[295]}}\)。

关于火星上存在智慧火星人的观念在 19 世纪末广泛传播。夏帕雷利（Schiaparelli）的“canali”（沟渠）观测，加上珀西瓦尔·洛厄尔（Percival Lowell）关于该主题的著作，共同塑造了当时的主流观点：火星是一颗正在枯竭、冷却、走向死亡的世界，拥有古文明并建造灌溉工程 \(^{\text{[296]}}\)。许多著名人物的观测与宣称进一步推动了所谓的“火星热”（Mars Fever）\(^{\text{[297]}}\)。如今高分辨率的火星地表制图已未发现任何居住痕迹，但关于火星智慧生命的伪科学推测仍在持续。这些推测类似于当年的“运河”幻象，基于航天器图像中的小尺度结构，例如所谓的“金字塔”与“火星人脸”（Face on Mars）\(^{\text{[298]}}\)。行星天文学家卡尔·萨根（Carl Sagan）在其著作《宇宙》（Cosmos）中写道：“火星已成为我们投射人类希望与恐惧的神话舞台。”\(^{\text{[238]}}\)

火星在文学作品中的描绘因其鲜明的红色外观与 19 世纪科学界对其潜在生命、乃至智慧生命的猜测而受到巨大刺激 \(^{\text{[299]}}\)。这激发了大量科幻作品，例如 H. G. 威尔斯的《世界大战》，描绘火星人为逃离垂死的故乡而入侵地球；雷·布拉德伯里的《火星编年史》，其中人类探险者意外毁灭了火星文明；以及埃德加·赖斯·巴勒斯（Edgar Rice Burroughs）的《巴索姆》系列、C. S. 刘易斯 1938 年的小说《沉寂星球》（Out of the Silent Planet）\(^{\text{[300]}}\)，和罗伯特·海因莱因（Robert A. Heinlein） 1960 年代之前的多篇作品 \(^{\text{[301]}}\)。自此之后，火星人的形象也延伸至动画。例如火星智能生物的滑稽形象“火星马文”（Marvin the Martian）首次于 1948 年的《Haredevil Hare》中出现，并在华纳兄弟的《乐一通》（Looney Tunes）动画中成为持续至今的文化符号 \(^{\text{[302]}}\)。在水手号与海盗号探测器传回火星无生命、无“运河”的图像后，这些旧有观念被抛弃；对许多科幻作家而言，这些新发现初看似限制，但最终“后海盗号时代”的火星知识本身又成为灵感源泉，例如金·斯坦利·罗宾逊（Kim Stanley Robinson）的《火星三部曲》 \(^{\text{[303]}}\)。
\subsection{参见} 
\begin{itemize}
\item 火星上的天文学（Astronomy on Mars）  
\item 火星纲要——火星的概览与主题指南（Outline of Mars – Overview of and topical guide to Mars）  
\item 火星探测任务列表（List of missions to Mars）  
\item 火星的磁场——火星过去的行星磁场（Magnetic field of Mars – Past magnetic field of the planet Mars）  
\item 火星矿物学（Mineralogy of Mars）
\end{itemize}
\subsection{注释}\\  
a.光学滤光片的波长分别为 635 nm、546 nm 和 437 nm，分别大致对应红、绿、蓝光。\\ 
b.最佳拟合椭球体（Best-fit ellipsoid）\\
\subsection{参考文献}

\begin{enumerate}
\item Simon J，Bretagnon P，Chapront J等 (1994年2月)。“月球和行星的进动公式和平均元素的数值表达式”。天文学和天体物理学。282(2):663-683。Bibcode:1994A & A。282 .. 663S。
\item 威廉姆斯D (2018)。"火星概况"。美国宇航局戈达德太空飞行中心。存档自原来的2020年3月17日。已检索3月22日2020。；平均异常 (deg) 19.412 = (平均经度 (deg) 355.45332) - (近日点经度 (deg) 336.04084)公共领域本文结合了此来源的文本，该文本位于公共领域。
\item 艾伦·CW，考克斯·安 (2000)。艾伦的天体物理量。施普林格科学与商业媒体。第294页。ISBN 978-0-387-95189-8。已存档从原始的2024年3月1日。已检索5月18日2022年。
\item Souami D，Souchay J (2012年7月)。“太阳系的不变平面”。天文学与天体物理学。543: 11。Bibcode:2012A & A...543A.133S。doi:10.1051/0004-6361/201219011。A133.
\item “2022年近日点的地平线批处理电话”(当rdot从负翻转为正时发生近日点)。太阳系动力学小组，喷气推进实验室。已存档从原始的2021年9月8日。已检索9月7日2021。
\item Seidelmann PK，Archinal BA，A'Hearn MF等人 (2007年)。“IAU/IAG制图坐标和旋转要素工作组的报告: 2006年”。天体力学与动力天文学。98(3):155-180。Bibcode:2007CeMDA .. 98 .. 155S。doi:10.1007/s10569-007-9072-y。
\item Grego P (2012年6月6日)。火星以及如何观察它。施普林格科学 + 商业媒体。p。 3。ISBN 978-1-4614-2302-7-通过互联网档案馆。
\item Lodders K,Fegley B (1998)。行星科学家的同伴。牛津大学出版社.p. 190。ISBN 978-0-19-511694-6。
\item Konopliv AS，Asmar SW，Folkner WM等人 (2011年1月)。“来自MRO、火星季节性重力和其他动力学参数的火星高分辨率重力场”。伊卡洛斯。211(1):401-428。Bibcode:2011Icar .. 211 .. 401K。doi:10.1016/j.icarus.2010.10.004。
\item Hirt C，Claessens SJ，Kuhn M等人 (2012年7月)。“火星的公里分辨率重力场: MGM2011”(PDF)。行星与空间科学。67(1):147-154.Bibcode:2012P & SS...67 .. 147小时。doi:10.1016/j.pss.2012.02.006。高密度脂蛋白:20.500.11937/32270。ISSN0032-0633。已存档从原始的2024年3月1日。已检索8月25日2019。  
\item Jackson AP，Gabriel TS，Asphaug EI (2018年3月1日)。“对太阳系巨型撞击器撞击前轨道的限制”。皇家天文学会月刊。474(3):2924-2936。arXiv:1711.05285。doi:10.1093/mnras/stx2901。ISSN0035-8711。已存档从原始的2022年4月23日。已检索4月23日2022年。 
\item 艾莉森·M，施蒙克·R。“Mars24太阳时钟-火星上的时间”。NASA GISS。已存档从2017年1月24日开始。已检索8月19日2021。
\item Archinal BA，Acton CH，A'Hearn MF等人 (2018年)。“IAU制图坐标和旋转要素工作组的报告: 2015年”。天体力学与动力天文学。130(3): 22。Bibcode:2018 CeMDA.130. .. 22A。doi:10.1007/s10569-017-9805-5。ISSN0923-2958。 
\item Mallama, A.(2007)。“火星的大小和反照率”。伊卡洛斯。192(2):404-416.Bibcode:2007年Icar .. 192 .. 404米。doi:10.1016/j.icarus.2007.07.011。
\item “大气和行星温度”。美国化学学会。2013年7月18日存档自原来的2023年1月27日。已检索1月3日2023。
\item “火星探测漫游者任务: 聚光灯”。Marsrover.nasa.gov。2007年6月12日。存档自原来的2013年11月2日。已检索8月14日2012。 公共领域本文结合了此来源的文本，该文本位于公共领域。
\item 夏普·T，戈登·J，蒂尔曼·N (2022年1月31日)。“火星的温度是多少？”。Space.com。已存档从原件到2020年4月22日。已检索3月14日2022年。
\item Hassler DM、Zeitlin C、wimmer-schweingruber RF等人 (2014年1月24日)。“用火星科学实验室的好奇号火星车测量火星表面辐射环境”。科学。343(6169) 1244797。表1和2。Bibcode:2014Sci...343D.386H。doi:10.1126/science.1244797。高密度脂蛋白:1874/309142。PMID24324275。S2CID33661472。  
\item Mallama A，希尔顿JL (2018年10月)。“计算天文历书的视在行星星等”。天文学与计算。25:10-24.arXiv:1808.01973。Bibcode:2018A & C .... 25...10米。doi:10.1016/j.ascom.2018.08.002。S2CID69912809。 
\item “百科全书-最聪明的身体”。IMCCE。已存档从2023年7月24日开始。已检索5月29日2023。
\item 巴洛NG(2008)。火星: 介绍其内部，表面和大气。剑桥行星科学。卷。8。剑桥大学出版社。第21页。ISBN 978-0-521-85226-5。
\item Rees MJ，ed. (2012年10月)。宇宙: 权威的视觉指南。纽约: Dorling Kindersley。pp. 160-161.ISBN 978-0-7566-9841-6。
\item “赤铁矿的诱惑”。科学 @ NASA。NASA.2001年3月28日。存档自原来的2010年1月14日。已检索12月24日2009年。
\item “火星: 事实”。NASA科学。2017年11月10日。已检索4月27日2025。
\item 哈利迪，A。N.,W ä nke，H.，Birck，J.-L.等人 (2001年)。“火星的生长、组成和早期分化”。空间科学评论。96(1/4):197-230。Bibcode:2001SSRv...96 .. 197小时。doi:10.1023/A:1011997206080。S2CID55559040。 
\item Zharkov VN (1993)。“木星在行星形成中的作用”。地球和行星的演化。华盛顿特区美国地球物理联合会地球物理专著系列。第74卷。pp。 7-17.Bibcode:1993GMS....74....7Z。doi:10.1029/GM074p0007。ISBN 978-1-118-66669-2。
\item Lunine JI钱伯斯，Morbidelli A等 (2003年)。“火星上的水的起源”。伊卡洛斯。165(1):1-8.Bibcode:2003年Icar .. 165 .... 1L。doi:10.1016/S0019-1035(03)00172-6。
\item 巴洛NG(1988年10月5日至7日)。H、弗雷 (编)。火星早期的条件: 陨石坑记录的限制。MEVTV火星早期构造和火山演化研讨会。LPI技术报告89-04。马里兰州伊斯顿: 月球与行星研究所。第15页。Bibcode:1989eamd.工作... 15B。
\item Nesvorn ý D (2018年6月)。“早期太阳系的动态演化”。天文学和天体物理学年度评论。56:137-174.arXiv:1807.06647。Bibcode:2018ARA & A .. 56 .. 137N。doi:10.1146/annurev-astro-081817-052028。
\item 耶格尔A (2008年7月19日)。“撞击可能改变了火星”。ScienceNews.org。存档自原来的2012年9月14日。已检索8月12日2008。
\item Minkel JR (2008年6月26日)。“研究表明，巨型小行星将火星的一半夷为平地”。科学美国人。已存档从原来的2014年9月4日。已检索4月1日2022年。
\item Chang K (2008年6月26日)。“巨大的流星撞击解释了火星的形状，报告说”。纽约时报。已存档从原来的2017年7月1日。已检索6月27日2008。
\item Ć uk M，Minton DA，Pouplin JL等人 (2020年6月16日)。“来自Deimos轨道倾角的过去火星环的证据”。天体物理学杂志。896(2): l28。arXiv:2006.00645。Bibcode:2020ApJ...896L..28C。doi:10.3847/2041-8213/ab974f。ISSN2041-8213。 
\item 新闻工作人员 (2020年6月4日)。“研究人员发现新的证据表明，火星曾经有巨大的环 | Sci.News”。Sci.News: 突发科学新闻。已存档从2023年11月7日开始。已检索11月7日2023。
\item “古代火星有戒指吗？”。Earthsky | 更新你的宇宙和世界。2020年6月5日。已存档从2023年11月7日开始。已检索11月7日2023。
\item Tanaka KL (1986)。“火星的地层学”。地球物理研究杂志。91(B13):E139-E158。Bibcode:1986JGR....91E.139T。doi:10.1029/JB091iB13p0E139。已存档从原件到2021年12月16日。已检索7月17日2019。
\item 威廉·K·哈特曼，Neukum，Gerhard (2001)。“陨坑年表与火星的演化”。空间科学评论。96(1/4):165-194.Bibcode:2001SSRv...96 .. 165小时。doi:10.1023/A:1011945222010。S2CID7216371。 
\item “ESA科学与技术-火星时代”。sci.esa.int。已存档从2023年8月29日开始。已检索12月7日2021。
\item 卡尔·L·米切尔，威尔逊，莱昂内尔 (2003)。“火星: 最近的地质活动: 火星: 一个地质活跃的星球”。天文学和地球物理学。44(4):4.16-4.20。Bibcode:2003A & G .... 44d .. 16米。doi:10.1046/j.1468-4004.2003.44416.x。
\item Russell P (2008年3月3日)。“陷入行动: 北极陡坡上的雪崩”。HiRISE运营中心。已存档从2022年4月7日开始。已检索3月28日2022年。
\item “HiRISE在火星上遇到雪崩”。NASA喷气推进实验室 (JPL)。2020年8月12日。已存档从2022年3月1日开始。已检索3月28日2022年。
\item Peplow M (2004年5月6日)。“火星是如何生锈的”。性质: news040503-6。doi:10.1038/news040503-6。已存档从原始的2022年11月28日。已检索3月10日2007。
\item NASA-火星在一分钟: 火星真的是红色的吗？ 已存档2014年7月20日在回程机(成绩单 已存档2015年11月6日在回程机)公共领域本文结合了此来源的文本，该文本位于公共领域。
\item St ä hler SC，Khan A，Banerdt WB等人 (2021年7月23日)。“火星核心的地震探测”。科学。373(6553):443-448.Bibcode:2021Sci...373 .. 443S。doi:10.1126/科学.abi7730。高密度脂蛋白:20.500.11850/498074。PMID34437118。S2CID236179579。已存档从原始的2024年3月1日。已检索10月17日2021。   
\item Samuel H，Drilleau，m é lanie，Rivoldini，Attilio等人 (2023年10月)。“火星核心上方富含熔融硅酸盐层的地球物理证据”。性质。622(7984):712-717。Bibcode:2023年自然622. 712s。doi:10.1038/s41586-023-06601-8。高密度脂蛋白:20.500.11850/639623。ISSN1476-4687。PMC10600000。PMID37880437。    
\item Nimmo F，Tanaka K (2005年)。“火星的早期地壳演化”。地球与行星科学年度评论。33(1):133-161.Bibcode:2005AREPS..33..133N。doi:10.1146/annurev.earth.33.092203.122637。S2CID45843366。 
\item “深度 | 火星”。NASA太阳系探索。2017年11月10日。已存档从原件到2021年12月21日。已检索1月15日2022年。
\item Kim D，Duran C，Giardini D等人 (2023年6月)。“围绕火星运行的表面波揭示的全球地壳厚度”。地球物理研究快报。50(12) e2023GL103482。Bibcode:2023年乔治。。5003482k。doi:10.1029/2023GL103482。高密度脂蛋白:20.500.11850/621318。
\item Wieczorek MA，Broquet A，McLennan SM等人 (2022年5月)。“洞察力对火星地壳全球特征的限制”。地球物理研究杂志: 行星。127(5) e2022JE007298。Bibcode:2022JGRE..12707298W。doi:10.1029/2022JE007298。高密度脂蛋白:10919/110830。
\item Huang Y，Chubakov V，Mantovani F等人 (2013年6月)。“生热元素和相关地中微通量的参考地球模型”。地球化学，地球物理学，地球系统。14(6):2003-2029年。arXiv:1301.0365。Bibcode:2013GGG....14.2003H。doi:10.1002/ggge.20129。
\item Fernando B，Daubar IJ，Charalambous C等人 (2023年10月)。“洞察力观测到的最大的马斯奎克的构造起源”。地球物理研究快报。50(20) e2023GL103619。Bibcode:2023乔治 .. 5003619F。doi:10.1029/2023GL103619。高密度脂蛋白:20.500.11850/639018。
\item Golombek M，Warner NH，Grant JA等人 (2020年2月24日)。“洞察号火星着陆点的地质”。自然地球科学。11(1014): 1014。Bibcode:2020年纳科。。11.1014克。doi:10.1038/s41467-020-14679-1。PMC7039939。PMID32094337。   
\item Banerdt WB，Smrekar SE，Banfield D等人 (2020年)。“洞察号火星任务的初步结果”。自然地球科学。13(3):183-189.Bibcode:2020年自然 .. 13 .. 183B。doi:10.1038/s41561-020-0544-y。
\item Khan A，Ceylan S，van Driel M等人 (2021年7月23日)。“来自InSight地震数据的火星上地幔结构”(PDF)。科学。373(6553):434-438。Bibcode:2021Sci...373 .. 434K。doi:10.1126/科学.abf2966。PMID34437116。S2CID236179554。已存档(PDF)从2022年1月20日开始。已检索11月27日2021。   
\item 汗A，黄，D.，Dur á n，C.等人 (2023年10月)。“火星核心顶部有液态硅酸盐层的证据”。性质。622(7984):718-723.Bibcode:2023年自然622. 718k。doi:10.1038/s41586-023-06586-4。高密度脂蛋白:20.500.11850/639367。PMC10600012。PMID37880439。   
\item Charalambous C，Pike WT，Kim D等人 (2025年8月28日)。“高度异质的火星地幔的地震证据”。科学。389(6763):899-903。Bibcode:2025Sci...389 .. 899C。doi:10.1126/科学.adk4292。ISSN0036-8075。PMID40875851。已检索9月4日2025。   
\item NASA Marsquake数据揭示了红色星球内部的凹凸不平的性质。NASA喷气推进实验室 (JPL)。已检索9月8日2025。 公共领域本文结合了此来源的文本，该文本位于公共领域。
\item Le Maistre S，Rivoldini A，Caldiero A等 (2023年6月14日)。“来自InSight无线电跟踪的火星的自旋状态和深层内部结构”。性质。619(7971):733-737。Bibcode:2023 Natur.619..733L。doi:10.1038/s41586-023-06150-0。高密度脂蛋白:2078.1/296670。ISSN1476-4687。PMID37316663。S2CID259162975。已存档从2023年10月12日开始。已检索7月3日2023。   
\item 雷恩E (2023年7月2日)。“火星有液体的内脏和奇怪的内部，洞察力表明”。Ars Technica。已存档从2023年7月3日开始。已检索7月3日2023。
\item Bi H，Sun D，Sun N，等人 (2025年9月3日)。“火星600公里实心内核的地震探测”。性质。645(8079)。施普林格科学与商业媒体:67-72。Bibcode:2025年自然645. .. 67B。doi:10.1038/s41586-025-09361-9。ISSN0028-0836。PMC12408336。PMID40903600。    
\item van der Lee S (2023年10月25日)。“火星很软”。性质。622(7984):699-700。Bibcode:2023年自然622. 699v。doi:10.1038/d41586-023-03151-x。PMID37880433。已存档从原件到2023年12月6日。已检索3月12日2024。 
\item Witze A (2023年10月25日)。“火星内部有一层令人惊讶的熔融岩石”。性质。623(7985): 20。Bibcode:2023年自然623. .. 20w。doi:10.1038/d41586-023-03271-4。PMID37880531。已存档从2024年1月9日开始。已检索3月12日2024。 
\item 谢文，德尚F，曹Y等人 (2024年)。“导热火星核心及其对发电机停止的影响”。科学进展。10(12) eadk1087。Bibcode:2024SciA...10K1087H。doi:10.1126/sciadv.adk1087。PMC10954222。PMID38507495。  
\item 麦克斯温·希，泰勒·GJ，怀亚特·MB (2009年5月)。“火星地壳的元素组成”。科学。324(5928):736-739。Bibcode:2009Sci...324 .. 736米。CiteSeerX10.1.1.654.4713。doi:10.1126/science.1165871。PMID19423810。S2CID12443584。   
\item Bandfield JL (2002年6月)。“火星上的全球矿物分布”。地球物理研究杂志: 行星。107(E6): 9-1-9-20。Bibcode:2002JGRE..107.5042B。CiteSeerX10.1.1.456.2934。doi:10.1029/2001JE001510。 
\item Christensen PR (2003年6月27日)。“火星表面的形态和组成: 火星奥德赛THEMIS结果”(PDF)。科学。300(5628):2056-2061。Bibcode:2003Sci...300.2056C。doi:10.1126/science.1080885。PMID12791998。S2CID25091239。已存档(PDF)从原件到2018年7月23日。已检索12月17日2018。   
\item 火星土壤 “可以支持生命”"。BBC新闻。2008年6月27日。已存档从2011年8月13日的原件。已检索8月7日2008。
\item Kounaves SP (2010)。“2007年凤凰号火星侦察员着陆器的湿化学实验: 数据分析和结果”。J.地球物理学。Res。115(E3):E00-E10。Bibcode:2009JGRE..114.0A19K。doi:10.1029/2008JE003084。S2CID39418301。  
\item Kounaves SP (2010)。“凤凰着陆点火星土壤中的可溶性硫酸盐”。地球物理研究快报。37(9) 2010GL042613: l09201。Bibcode:2010年乔治。。37.9201K。doi:10.1029/2010GL042613。S2CID12914422。 
\item David L (2013年6月13日)。有毒的火星: 宇航员必须处理红色星球上的高氯酸盐。Space.com。已存档从原件到2020年11月20日。已检索11月26日2018。
\item 样本I (2017年7月6日)。“火星上覆盖着有毒化学物质，可以消灭活生物体，测试显示”。守护者。已存档从原件到2021年2月18日。已检索11月26日2018。
\item Verba C (2009年7月2日)。"Dust Devil Etch-A-草图 (ESP_013751_1115)"。HiRISE运营中心。亚利桑那大学。已存档从2012年2月1日起。已检索3月30日2022年。
\item Schorghofer，Norbert，Aharonson，Oded，Khatiwala，Samar (2002)。“火星上的斜坡条纹: 与表面特性和水的潜在作用的相关性”(PDF)。地球物理研究快报。29(23):41-1.Bibcode:2002年乔治。。29.2126s。doi:10.1029/2002GL015889。已存档(PDF)从原件到2020年8月3日。已检索8月25日2019。
\item G á nti，Tibor (2003)。“黑暗的沙丘斑点: 火星上可能的生物标志物？”生命起源和生物圈的进化。33(4):515-557。Bibcode:2003OLEB...33..515G。doi:10.1023/A:1025705828948。PMID14604189。S2CID23727267。  
\item 威廉姆斯M (2016年11月21日)。“火星上的辐射有多严重？”。Phys.org。已存档从原始的2023年4月4日。已检索4月9日2023。
\item Wall M (2013年12月9日)。火星上的辐射 “可管理” 载人飞行任务，好奇号火星车揭示。Space.com。已存档从原件到2020年12月15日。已检索4月9日2023。
\item “火星辐射环境与国际空间站的比较”。NASA喷气推进实验室 (JPL)。2003年3月13日。已存档从2023年4月9日开始。已检索4月9日2023。
\item 巴黎A，戴维斯E，Tognetti L等人 (2020年4月27日)。“Hellas Planitia的预期熔岩管”。arXiv:2004.13156v1[astro-ph.EP]。
\item Valentine，Theresa，Amde，Lishan (2006年11月9日)。《磁场与火星》。火星全球测量师 @ NASA。已存档从原件到2012年9月14日。已检索7月17日2009年。 公共领域本文结合了此来源的文本，该文本位于公共领域。
\item 尼尔-琼斯N，奥卡罗尔C。“新地图提供了更多的证据火星曾经像地球”。NASA/戈达德太空飞行中心。存档自原来的2012年9月14日。已检索12月4日2011年。 公共领域本文结合了此来源的文本，该文本位于公共领域。
\item 什么火星地图通过时间得到了正确的 (和错误的)。国家地理。2016年10月19日存档自原来的2021年2月21日。已检索1月15日2022年。
\item “行星名称: 行星和卫星上命名特征的类别”。国际天文学联合会 (IAU) 行星系统命名工作组 (WGPSN)。USGS天体地质科学中心。已存档从2017年12月5日的原件。已检索4月18日2022年。
\item “维京和火星的资源”(PDF)。人类到火星: 五十年的任务规划，1950-2000。存档自原来的(PDF)2019年7月14日。已检索3月10日2007。 公共领域本文结合了此来源的文本，该文本位于公共领域。
\item Tanaka KL，Coles KS，Christensen PR，eds。(2019年)，"Syrtis Major (MC-13)",火星地图集: 绘制其地理和地质,剑桥: 剑桥大学出版社，pp。 136-139,doi:10.1017/9781139567428.018,ISBN 978-1-139-56742-8,S2CID 240843698,已存档从2024年3月1日起,检索到1月18日2022年
\item "极坐标"。亚利桑那州立大学的火星教育。已存档从原件到2021年5月28日。已检索1月6日2022年。
\item Davies ME，Berg RA (1971年1月10日)。“火星的初步控制网”。地球物理研究杂志。76(2):373-393。Bibcode:1971JGR....76..373D。doi:10.1029/JB076i002p00373。已存档从2024年3月1日起。已检索3月22日2022年。
\item 阿基纳，B。A.,卡普林格，M.(2002年秋季)。“火星、子午线和梅特: 探索火星经度”。美国地球物理联合会，2002年秋季会议。22: P22D-06。Bibcode:2002AGUFM.P22D..06A。
\item de Vaucouleurs G,戴维斯我,Sturms FM Jr (1973)，“水手9区域坐标系统”，地球物理研究杂志,78(20):4395-4404,Bibcode:1973JGR....78.4395D,doi:10.1029/JB078i020p04395
\item NASA (2007年4月19日)。“火星全球测量师: MOLA MEGDRs”。geo.pds.nasa.gov。存档自原来的2011年11月13日。已检索6月24日2011年。
\item Ardalan AA，Karimi R，Grafarend EW (2009)。“一个新的参考等势面，以及火星的参考椭球”。地球、月球和行星。106(1):1-13.doi:10.1007/s11038-009-9342-7。ISSN0167-9295。S2CID119952798。  
\item Zeitler, W.,Ohlhof, T.,埃伯纳，H. (2000)。“全球火星控制点网络的重新计算”。摄影测量工程与遥感。66(2):155-161.CiteSeerX10.1.1.372.5691。 
\item Lunine CJ (1999)。地球: 可居住世界的演变。剑桥大学出版社.p. 183。ISBN 978-0-521-64423-5。
\item “ESA科学与技术-使用iMars: 查看MC11四边形的Mars Express数据”。sci.esa.int。已存档从原件到2021年12月29日。已检索12月29日2021。
\item 张凯 (2023年4月15日)。新的火星地图让你 “一次看到整个星球” -- 科学家从一个阿联酋轨道器收集了3000张图像，创建了这颗红色星球迄今为止最漂亮的地图集。。纽约时报。已存档从2023年5月16日开始。已检索4月15日2023。
\item 工作人员 (2023年4月16日)。欢迎来到火星!加州理工学院令人瞠目结舌的5.7 Terapixel穿越红色星球的虚拟探险“。科技。已存档从原始的2023年4月16日。已检索4月6日2023。
\item “火星地图集: 奥林巴斯蒙斯”。NASA的火星探测计划。已存档从2023年3月29日开始。已检索3月30日2022年。
\item Plescia JB (2004年)。“火星火山的形态计量特性”。J.地球物理学。Res。109(E3) 2002JE002031: E03003.Bibcode:2004JGRE..109.3003P。doi:10.1029/2002JE002031。
\item 科明斯NF (2012)。发现本质的宇宙。W。H、弗里曼。第148页。ISBN 978-1-4292-5519-6。
\item 申克·P (2012)。“位于灶神星南极的地质上最近的巨型撞击盆地”。科学。336(6082):694-697。Bibcode:2012Sci...336 .. 694S。doi:10.1126/science.1223272。PMID22582256。S2CID206541950。  
\item 安德鲁斯-汉娜JC，Zuber MT,Banerdt WB (2008)。“北欧盆地与火星地壳二分法的起源”。性质。453(7199):1212-1215。Bibcode:2008 Natur.453.1212A。doi:10.1038/nature07011。PMID18580944。S2CID1981671。  
\item 崔C (2021年10月1日)。“火星: 我们对红色星球的了解”。Space.com。已存档从2022年1月6日开始。已检索1月6日2022年。
\item Moskowitz C (2008年6月25日)。“巨大的影响造成了火星的人格分裂”。Space.com。已存档从2022年1月6日开始。已检索1月6日2022年。
\item Wright S (2003年4月4日)。“地球和火星上小型撞击坑的红外分析”。匹兹堡大学。存档自原来的2007年6月12日。已检索2月26日2007。
\item 沃格特GL (2008)。火星的风景(PDF)。纽约: 施普林格。第44页。doi:10.1007/978-0-387-75468-0。ISBN  978-0-387-75467-3。已存档从2024年3月1日起。已检索3月31日2022年。
\item “欧空局科学与技术-海拉斯盆地内的火山口”。sci.esa.int。已存档从2022年1月2日开始。已检索1月2日2022年。
\item 罗德里格CM。《火星的地理》。Home.csulb.edu。已存档从2022年1月30日开始。已检索2月20日2022年。
\item 第41届月球与行星科学会议 (2010年)(PDF)。已存档(PDF)从2022年1月30日开始。已检索1月31日2022年。
\item Wetherill GW (1999)。“与估计对火星和月球的相对影响率相关的问题”。地球、月球和行星。9(1-2):227-231.Bibcode:1974年月亮。9。。227W。doi:10.1007/BF00565406。S2CID120233258。 
\item Costard FM (1989)。“火星水岩石圈中挥发物的空间分布”。地球、月球和行星。45(3):265-290.Bibcode:1989EM & P...45 .. 265C。doi:10.1007/BF00057747。S2CID120662027。 
\item 萨根C福克斯·P (1975年8月)。“火星运河: 水手9号之后的评估”。伊卡洛斯。25(4):602-612。Bibcode:1975年Icar。25 .. 602S。doi:10.1016/0019-1035(75)90042-1。已检索3月22日2022年。
\item Wolpert S (2012年8月9日)。“加州大学洛杉矶分校的科学家在火星上发现了板块构造”。加州大学洛杉矶分校。存档自原来的2012年8月12日。已检索8月13日2012。
\item 林A (2012年6月4日)。“Valles Marineris断层带的结构分析: 火星上大规模走滑断层的可能证据”。岩石圈。4(4):286-330。Bibcode:2012Lsphe...4 .. 286Y。doi:10.1130/L192.1。
\item 库欣，G.E.,泰特斯，T.N.,Wynne, J.J.等人 (2007年)。“Themis观察到火星上可能的洞穴天窗”(PDF)。月球和行星科学XXXVIII。已存档(PDF)从2011年9月15日的原件。已检索8月2日2007。
\item NAU研究人员在火星上发现了可能的洞穴。NAU内部。第4卷，第12号。北亚利桑那大学。2007年3月28日。已存档从2007年8月28日的原件。已检索5月28日2007。
\item Rincon P(2007年3月17日)。"火星上发现的 “洞穴入口”。BBC新闻。已存档从2009年9月30日的原件。已检索5月28日2007。
\item “火星的洞穴 | 美国地质调查局”。美国地质勘探局。已存档从原件到2021年12月31日。已检索1月12日2022年。
\item Portyankina G (2014)。“Araneiform”。行星地貌大全:。p. 1.doi:10.1007/978-1-4614-9213-9_540-1。ISBN 978-1-4614-9213-9。
\item Cantor，B (2007)。“MOC对2001年环绕火星的沙尘暴的观测”。伊卡洛斯。186(1):60-96。Bibcode:2007Icar .. 186...60摄氏度。doi:10.1016/j.icarus.2006.08.019。
\item Claudin，P，Andreotti B (2006)。“火星，金星，地球上的风沙沙丘和水下涟漪的缩放定律”。地球和行星科学快报。252(1-2):30-44。arXiv:cond-mat/0603656。Bibcode:2006E & PSL.252. .. 30C。doi:10.1016/j.epsl.2006.09.004。S2CID13910286。 
\item Philips T (2001年1月31日)。“火星上的太阳风”。科学 @ NASA。存档自原来的2011年8月18日。已检索4月22日2022年。 公共领域本文结合了此来源的文本，该文本位于公共领域。
\item 格罗斯曼L (2011年1月20日)。多次小行星撞击可能杀死了火星的磁场。有线。已存档从原件到2013年12月30日。已检索3月30日2022年。
\item Jakosky BM (2022年4月1日)。火星是如何失去大气层和水的？。今日物理。75(4):62-63.Bibcode:2022PhT。75d .. 62J。doi:10.1063/PT.3.4988。ISSN0031-9228。S2CID247882540。  
\item 伦丁，R (2004)。太阳风引起的火星大气侵蚀: 火星快车ASPERA-3的第一个结果。科学。305(5692):1933年-1936年。Bibcode:2004Sci...305.1933L。doi:10.1126/science.1101860。PMID15448263。S2CID28142296。  
\item 博隆金AA (2009)。火星上的人工环境。柏林海德堡: 施普林格。pp. 599-625。ISBN 978-3-642-03629-3。
\item 南希阿特金森 (2007年7月17日)。“火星着陆方法: 将大型有效载荷运送到红色星球的表面”。今天的宇宙。已存档从2010年8月5日的原件。已检索9月18日2007。
\item 卡尔·MH (2006)。火星表面。卷6。剑桥大学出版社。第16页。ISBN 978-0-521-87201-0。
\item “火星事实 | 关于火星的一切”。NASA的火星探测计划。2017年11月10日。已存档从2023年10月10日开始。已检索12月27日2021。
\item Mahaffy PR (2013年7月19日)。“好奇号火星车火星大气中气体的丰度和同位素组成”。科学。341(6143):263-266。Bibcode:2013Sci...341 .. 263米。doi:10.1126/science.1237966。PMID23869014。S2CID206548973。  
\item Lemmon MT (2004)。“火星漫游者的大气成像结果”。科学。306(5702):1753-1756.Bibcode:2004Sci。306.1753l。doi:10.1126/science.1104474。PMID15576613。S2CID5645412。  
\item 样本I (2018年6月7日)。Nasa火星探测器在古代湖床中发现有机物。守护者。已存档从原件到2018年6月18日。已检索6月12日2018。
\item Mumma MJ(2009年2月20日)。“2003年夏季火星上甲烷的强烈释放”(PDF)。科学。323(5917):1041-1045。Bibcode:2009Sci。323.1041米。doi:10.1126/science.1165243。PMID19150811。S2CID25083438。已存档(PDF)从2012年3月13日的原件。已检索11月1日2009年。   
\item 弗兰克·L，忘记F(2009年8月6日)。“火星上观察到的甲烷变化无法用已知的大气化学和物理学解释”。性质。460(7256):720-723.Bibcode:2009 Natur.460..720L。doi:10.1038/nature08228。PMID19661912。S2CID4355576。  
\item 奥兹，C.，夏尔马，M.(2005)。“有橄榄石，会产生气体: 蛇纹石化和火星上甲烷的非生物成因”。地球物理研究快报。32(10): l10203。Bibcode:2005乔治。。3210203O。doi:10.1029/2005GL022691。S2CID28981740。 
\item Webster CR，Mahaffy PR，Pla-Garcia J等人 (2021年6月)。“火星甲烷的昼夜差异表明Gale crater的夜间收容”。天文学与天体物理学。650: a166。Bibcode:2021A & A...650A.166W。doi:10.1051/0004-6361/202040030。ISSN0004-6361。S2CID236365559。  
\item Jones N，Steigerwald B，Brown D等人 (2014年10月14日)。“NASA任务首次看到火星高层大气”。NASA.已存档从原件到2014年10月19日。已检索10月15日2014。 公共领域本文结合了此来源的文本，该文本位于公共领域。
\item 赖特·K (2022年3月22日)。“在火星上测量的声速”。物理。1543。Bibcode:2022年PhyOJ .. 15...43W。doi:10.1103/物理15.43。S2CID247720720。已存档从原始的2022年4月12日。已检索4月6日2022年。 
\item 莫里斯·S，奇德·B，默多克·N等人 (2022年4月1日)。“火星音景的原位记录”。性质。605(7911):653-658。Bibcode:2022年自然605.653m。doi:10.1038/s41586-022-04679-0。ISSN0028-0836。PMC9132769。PMID35364602。S2CID247865804。    
\item 周D (2021年12月7日)。“在紫外线辉光中，阿联酋轨道飞行器发现了火星上的极光”。NBC新闻。已存档从2021年12月7日开始。已检索12月7日2021。
\item “火星上的极光-NASA科学”。science.nasa.gov。2015年5月10日已存档从原件到2015年5月14日。已检索5月12日2015。 公共领域本文结合了此来源的文本，该文本位于公共领域。
\item Brown D，neal-jones N，Steigerwald B等人 (2015年3月18日)。NASA航天器探测到火星周围的极光和神秘的尘埃云。NASA.释放15-045。已存档从原件到2015年3月19日。已检索3月18日2015。 公共领域本文结合了此来源的文本，该文本位于公共领域。
\item Deighan J，Jain SK，Chaffin MS等人 (2018年10月)。“在火星上发现质子极光”。自然天文学。2(10):802-807.Bibcode:2018NatAs...2 .. 802D。doi:10.1038/s41550-018-0538-5。ISSN2397-3366。S2CID105560692。已存档从2022年5月22日开始。已检索4月5日2022年。  
\item 施耐德NM、Jain SK、Deighan J等人 (2018年8月16日)。“2017年9月空天天气事件期间火星上的全球极光”。地球物理研究快报。45(15):7391-7398。Bibcode:2018年乔治。。45.7391s。doi:10.1029/2018GL077772。高密度脂蛋白:10150/631256。S2CID115149852。 
\item Webster G，neal-jones N，Scott J等人 (2017年9月29日)。“大型太阳风暴引发全球极光，并使火星表面的辐射水平翻倍”。NASA.存档自原来的2017年10月1日。已检索1月9日2018。 公共领域本文结合了此来源的文本，该文本位于公共领域。
\item 古德曼JC (1997年9月22日)。“火星气候的过去，现在和可能的未来”。麻省理工学院。存档自原来的2010年11月10日。已检索2月26日2007。
\item "火星的沙漠表面..."MGCM新闻稿。NASA.存档自原来的2007年7月7日。已检索2月25日2007。 公共领域本文结合了此来源的文本，该文本位于公共领域。
\item 克鲁格J (1992年9月1日)。“火星，地球的形象”。发现杂志。13(9): 70。Bibcode:1992光盘... 13... 70k。已存档从2012年4月27日的原件。已检索11月3日2009年。
\item Hille K (2015年9月18日)。“火星沙尘暴的事实和虚构”。NASA。已存档从2016年3月2日的原件。已检索12月25日2021。
\item Philips T (2001年7月16日)。“吞噬地球的沙尘暴”。科学 @ NASA。存档自原来的2006年6月13日。已检索6月7日2006。 公共领域本文结合了此来源的文本，该文本位于公共领域。
\item “火星上的干冰”。NASA科学。2011年4月6日。已检索5月1日2025。
\item “火星北极火山口的水冰”。欧空局。2005年7月28日。存档自原来的2010年2月9日。已检索3月19日2010。
\item Whitehouse D (2004年1月24日)。“水和火星的悠久历史”。BBC新闻。已存档从2009年1月11日的原件。已检索3月20日2010。
\item Holt JW，Safaeinili A，Plaut JJ，et al. (2008年11月21日)。“火星南部中纬度地区埋藏冰川的雷达探测证据”。科学。322(5905):1235-1238。Bibcode:2008Sci...322.1235H。doi:10.1126/science.1164246。高密度脂蛋白:11573/67950。ISSN0036-8075。JSTOR20145331。PMID19023078。S2CID36614186。已存档从2022年4月22日开始。已检索4月22日2022年。    
\item Byrne, Shane, Ingersoll, Andrew P.(2003)。“火星南极冰特征的升华模型”。科学。299(5609):1051-1053。Bibcode:2003Sci...299.1051B。doi:10.1126/science.1080148。PMID12586939。S2CID7819614。  
\item “火星南极冰层深而宽”。NASA.2007年3月15日。存档自原来的2009年4月20日。已检索3月16日2007。 公共领域本文结合了此来源的文本，该文本位于公共领域。
\item “NASA - NASA漫游者发现火星大气变化的线索”。NASA.存档自原来的2018年12月26日。已检索10月19日2014。 公共领域本文结合了此来源的文本，该文本位于公共领域。
\item Holdmann JL (2005年5月7日)。“在当前火星环境条件下通过液态水流动形成火星沟壑”(PDF)。地球物理研究杂志。110(E5) 2004JE002261: Eo5004.Bibcode:2005JGRE..110.5004H。CiteSeerX10.1.1.596.4087。doi:10.1029/2004JE002261。高密度脂蛋白:2060/20050169988。S2CID1578727。存档自原来的(PDF)2008年10月1日。已检索9月17日2008。    “像现在这样的条件发生在火星上，在液态水的温度-压力稳定状态之外。“液态水在地球上的最低海拔和低纬度地区通常是稳定的，因为大气压力大于蒸汽压力赤道地区的水和地表温度在一天中的部分时间可以达到273 K [Haberle等。.,2001]'
\item “火星上的二氧化碳降雪”。NASA喷气推进实验室 (JPL)。2012年9月11日。已检索4月19日2025。
\item Kerr RA (2005年3月4日)。“火星上的冰还是熔岩海？一场跨大西洋的辩论爆发了 ”。科学。307(5714):1390-1391。doi:10.1126/science.307.5714.1390a。PMID15746395。S2CID38239541。  
\item Jaeger WL (2007年9月21日)。“阿萨巴斯卡山谷，火星: 一个熔岩覆盖的通道系统”。科学。317(5845):1709-1711.Bibcode:2007Sci...317.1709J。doi:10.1126/science.1143315。PMID17885126。S2CID128890460。  
\item 卢奇塔，B.K.,罗萨诺瓦，C。E.(2003年8月26日)。“Valles Marineris; 火星大峡谷”。美国地质勘探局。存档自原来的2011年6月11日。已检索3月11日2007。 公共领域本文结合了此来源的文本，该文本位于公共领域。
\item Murray JB (2005年3月17日)。“来自火星快车高分辨率立体相机的证据，证明了火星赤道附近的冰冻海洋”。性质。434(703):352-356。Bibcode:2005自然434..352米。doi:10.1038/nature03379。PMID15772653。S2CID4373323。  
\item 克拉多克，r.a.，Howard，a.d. (2002)。“在温暖潮湿的早期火星上降雨的情况”。地球物理研究杂志。107(E11):21-1.Bibcode:2002JGRE..107.5111C。CiteSeerX10.1.1.485.7566。doi:10.1029/2001JE001505。 
\item Malin MC，Edgett KS (2000年6月30日)。“火星上最近地下水渗漏和地表径流的证据”。科学。288(5475):2330-2335。Bibcode:2000Sci...288.2330米。doi:10.1126/science.288.5475.2330。PMID10875910。S2CID14232446。  
\item “NASA的图像表明，水仍然在火星上短暂喷涌而出”。NASA.2006年12月6日。存档自原来的2011年8月7日。已检索12月6日2006。 公共领域本文结合了此来源的文本，该文本位于公共领域。
\item “水最近在火星上流动”。英国广播公司。2006年12月6日。已存档从2011年8月30日的原件。已检索12月6日2006。
\item “水可能仍在火星上流动，NASA照片建议”。NASA.2006年12月6日。已存档从2007年1月2日的原件。已检索4月30日2006。 公共领域本文结合了此来源的文本，该文本位于公共领域。
\item 刘易斯，K.W.，Aharonson, O.(2006)。“使用立体图像对Eberswalde火山口的分流扇进行地层分析”(PDF)。地球物理研究杂志。111(E06001): E06001.Bibcode:2006JGRE..111.6001L。doi:10.1029/2005JE002558。已存档(PDF)从原件到2020年8月3日。已检索8月25日2019。
\item 松原，Y.，霍华德，A.D.，德拉蒙德，S.A.(2011)。“火星早期的水文学: 湖泊盆地”。地球物理研究杂志。116(E04001): E04001。Bibcode:2011JGRE..116.4001M。doi:10.1029/2010JE003739。
\item “火星 'Berries' 中的矿物质增加了水的故事”(新闻稿)。NASA.2004年3月3日。存档自原来的2007年11月9日。已检索6月13日2006。 公共领域本文结合了此来源的文本，该文本位于公共领域。
\item “火星探测漫游者任务: 科学”。NASA.2007年7月12日。存档自原来的2010年5月28日。已检索1月10日2010。 公共领域本文结合了此来源的文本，该文本位于公共领域。
\item Elwood Madden ME，Bodnar RJ，Rimstidt JD (2004年10月)。“黄金钠岩是火星上受水限制的化学风化的指标”。性质。431(7010):821-823.doi:10.1038/nature02971。ISSN0028-0836。PMID15483605。S2CID10965423。已存档从2022年3月3日的原件。已检索4月22日2022年。   
\item “火星漫游者调查火星过去的迹象”。NASA喷气推进实验室 (JPL)。2007年12月10日。已存档从2022年7月6日开始。已检索4月5日2022年。
\item “NASA - NASA火星探测器发现水沉积的矿脉”。NASA.2011年12月7日。存档自原来的2017年6月15日。已检索8月14日2012。 公共领域本文结合了此来源的文本，该文本位于公共领域。
\item Lovett RA (2011年12月8日)。“漫游者在火星早期发现了水的“ 防弹 ”证据”。国家地理。存档自原来的2021年5月1日。已检索3月31日2022年。
\item Lovett RA (2012年6月26日)。“火星内部有 '海洋' 的水吗？”。国家地理。存档自原来的2021年4月28日。已检索3月31日2022年。
\item McCubbin FM，Hauri EH，Elardo SM等人 (2012年8月)。“火星地幔的含水融化产生了贫化和富集的雪尔戈蒂特”。地质学。40(8):683-686。Bibcode:2012Geo .... 40 .. 683米。doi:10.1130/G33242.1。ISSN1943-2682。 
\item 韦伯斯特G，布朗D (2013年3月18日)。“好奇号火星探测器看到了水存在的趋势”。NASA.存档自原来的2013年4月24日。已检索3月20日2013。 公共领域本文结合了此来源的文本，该文本位于公共领域。
\item Rincon P (2013年3月19日)。“好奇心打破岩石，露出耀眼的白色内饰”。BBC新闻。英国广播公司。已存档从原件到2021年3月8日。已检索3月19日2013。
\item “NASA证实了液态水在今天的火星上流动的证据”。NASA.2015年9月28日已存档从原来的2015年9月28日。已检索9月28日2015。 公共领域本文结合了此来源的文本，该文本位于公共领域。
\item Drake N (2015年9月28日)。NASA在火星上发现了 “确定的” 液态水。国家地理新闻。存档自原来的2015年9月30日。已检索9月29日2015。
\item Ojha L，Wilhelm MB，Murchie SL等人 (2015年)。“火星上反复出现的斜坡线中的水合盐的光谱证据”。自然地球科学。8(11):829-832。Bibcode:2015年自然。8。829O。doi:10.1038/ngeo2546。S2CID59152931。 
\item 莫斯科维茨C.“今天火星上的水流，NASA宣布”。科学美国人。已存档从原件到2021年5月15日。已检索9月29日2015。
\item McEwen A、Lujendra O、Dundas C等人 (2011年8月5日)。“火星暖坡上的季节性水流”。科学。333(6043):740-743。Bibcode:2011Sci...333 .. 740米。doi:10.1126/science.1204816。PMID21817049。S2CID10460581。存档自原来的2015年9月29日。已检索9月28日2015。  
\item Dundas CM，McEwen AS，Chojnacki M等人 (2017年12月)。“火星上反复出现的斜坡线的颗粒流表明液态水的作用有限”。自然地球科学。10(12):903-907。Bibcode:2017NatGe .. 10 .. 903D。doi:10.1038/s41561-017-0012-5。高密度脂蛋白:10150/627918。ISSN1752-0908。S2CID24606098。已存档从2017年11月22日的原件。已检索4月3日2022年。  
\item Schorghofer N (2020年2月12日)。“火星: 巨石后番红花融化的定量评估”。天体物理学杂志。890(1): 49。Bibcode:2020年apj... 890... 49s。doi:10.3847/1538-4357/ab612f。ISSN1538-4357。S2CID213701664。  
\item Wray JJ (2021年5月30日)。“火星上的当代液态水？”。地球与行星科学年度评论。49(1):141-171。Bibcode:2021AREPS..49..141W。doi:10.1146/annurev-earth-072420-071823。ISSN0084-6597。S2CID229425641。已存档从原件到2021年5月3日。已检索4月3日2022年。  
\item 头，j.w.(1999)。“火星上可能的古代海洋: 来自火星轨道激光高度计数据的证据”。科学。286(5447):2134-7.Bibcode:1999Sci。286.2134H。doi:10.1126/science.286.5447.2134。PMID10591640。S2CID35233339。  
\item 考夫曼M (2015年3月5日)。“火星有海洋，科学家说，指向新的数据”。纽约时报。已存档从原件到2020年3月7日。已检索3月5日2015。
\item 样本I (2018年12月21日)。“火星快车传回了充满冰的科罗列夫陨石坑的图像”。守护者。已存档从2020年2月8日开始。已检索12月21日2018。
\item “EPA; 五大湖; 物理事实”。2010年10月29日。存档自原来的2010年10月29日。已检索2月15日2023。
\item “火星的冰沉积物拥有与苏必利尔湖一样多的水”。NASA.2016年11月22日已存档从原件到2018年12月26日。已检索11月23日2016。 公共领域本文结合了此来源的文本，该文本位于公共领域。
\item 工作人员 (2016年11月22日)。“扇形地形导致在火星上发现了埋藏的冰”。NASA.已存档从原件到2018年12月26日。已检索11月23日2016。 公共领域本文结合了此来源的文本，该文本位于公共领域。
\item Mitrofanov I，Malakhov A，Djachkova M，等人 (2022年3月)。“火星上Valles Marineris中部异常高的氢丰度的证据”。伊卡洛斯。374114805。Bibcode:2022年Icar .. 37414805米。doi:10.1016/j.icarus.2021.114805。S2CID244449654。 
\item Badescu V (2009年)。火星: 未来的能源和物质资源(插图版)。施普林格科学与商业媒体。第600页。ISBN 978-3-642-03629-3。已存档从2023年3月4日开始。已检索5月20日2016。
\item Petropoulos AE，Longuski JM，Bonfiglio EP (2000)。“从金星、地球和火星通过重力辅助到达木星的轨迹”。宇宙飞船和火箭杂志。37(6)。美国航空航天学会 (AIAA):776-783。Bibcode:2000JSpRo。。37 .. 776P。doi:10.2514/2.3650。ISSN0022-4650。 
\item 泰勒C (2020年7月9日)。“欢迎来到云城: 去金星的理由，而不是火星”。可混搭。已存档从2022年10月21日开始。已检索10月21日2022年。
\item Vitagliano A (2003年)。“火星随时间的轨道偏心率”。Solex。那不勒斯费德里科二世大学degli Studi。存档自原来的2007年9月7日。已检索7月20日2007。
\item Meeus J(2003年3月)。“火星最后一次这么近是什么时候？”。国际天文馆协会。存档自原来的2011年5月16日。已检索1月18日2008。
\item Laskar J (2003年8月14日)。“火星对立的底漆”。IMCCE，巴黎天文台。已存档从2011年11月13日的原件。已检索10月1日2010。 (Solex结果) 已存档2012年8月9日在回程机
\item Tillman NT，Dobrijevic D (2022年1月20日)。“到火星需要多长时间？”。Space.com。已检索8月23日2024。
\item “火星反对派”。NASA的火星探测计划。2017年11月10日。已存档从2023年10月5日开始。已检索12月7日2021。
\item “为什么火星时而明亮，时而暗淡？”。EarthSky。2021年10月5日。已存档从2021年12月7日开始。已检索12月7日2021。
\item “亲密接触: 反对派的火星”。欧空局/哈勃。2005年11月3日。已存档从原来的2015年9月10日。已检索4月1日2022年。
\item Mallama, A.(2011)。“行星量级”。天空和望远镜。121(1):51-56。Bibcode:2011S & T...121a .. 51米。
\item “火星接近”。NASA的火星探测计划。2017年11月10日。已存档从2019年11月8日的原件。已检索1月18日2022年。
\item “幻灯片2地球望远镜的火星视图”。红色星球: 火星调查。月球和行星研究所。已存档从原件到2021年5月18日。已检索11月28日2011年。
\item Zeilik M (2002年)。天文学: 进化中的宇宙(第9版)。剑桥大学出版社。第14页。ISBN 978-0-521-80090-7。
\item "火卫一近距离检查"。ESA网站。已存档从2012年1月14日开始。已检索6月13日2006。
\item “行星名称”。planetarynames.wr.usgs.gov。已存档从原件到2023年12月30日。已检索5月30日2022年。
\item "火卫一"。NASA太阳系探索。2019年12月19日。已存档从2022年1月12日开始。已检索1月12日2022年。
\item “解释火星卫星的诞生”。AAS新星。美国天文学会。2016年9月23日已存档从原件到2021年12月13日。已检索12月13日2021。
\item 阿德勒·M，欧文·W，里德尔·J (2012年6月)。使用MRO光学导航相机为火星样品返回做准备(PDF)。火星探索的概念和方法。2012年6月12-14日。德克萨斯州休斯顿。4337。Bibcode:2012LPICo1679.4337A。已存档从原件到2018年12月26日。已检索8月28日2012。
\item 安德特·TP，罗森布拉特，P.，P ä tzold，M.等人 (2010年5月7日)。“精确的质量测定和火卫一的性质”。地球物理研究快报。37(L09202): L09202。Bibcode:2010乔治。。37.9202a。doi:10.1029/2009GL041829。
\item Giuranna M, Roush, T.L.,Duxbury，T.等人 (2010)。Phobos的PFS/MEx和TES/MGS热红外光谱的组成解释(PDF)。欧洲行星科学大会摘要，第5卷。已存档(PDF)从2011年5月12日的原件。已检索10月1日2010。
\item Bagheri A，Khan A，Efroimsky M等人 (2021年2月22日)。“火卫一和火卫二是被破坏的共同祖先的残余物的动态证据”。自然天文学。5(6):539-543。Bibcode:2021年纳塔斯。。5 .. 539B。doi:10.1038/s41550-021-01306-2。S2CID233924981。 
\item Rabkin ES (2005)。火星: 人类想象力之旅。康涅狄格州韦斯特波特: pp. 9-11.ISBN 978-0-275-98719-0。
\item 汤普森·何 (1970)。麦卡尔: 贝丝山之神。莱顿: E。J.布里尔。第125页。
\item 诺瓦科维奇B (2008年)。“Senenmut: 古埃及天文学家”。贝尔格莱德天文台的出版物。85:19-23.arXiv:0801.1331。Bibcode:2008POBeo .. 85...19N。
\item 北JD(2008)。宇宙: 天文学和宇宙学的图解历史。芝加哥大学出版社。pp. 48-52。ISBN 978-0-226-59441-5。
\item Swerdlow NM (1998)。“Synodic现象的周期性和可变性”。巴比伦的行星理论。普林斯顿大学出版社。pp. 34-72.ISBN 978-0-691-01196-7。
\item 西塞罗 (1896)。De Natura Deorum[论神的本质]。由弗朗西斯·布鲁克斯翻译。伦敦: Methuen。
\item NASA (2022年10月9日)。“关于火星的一切”。mars.nasa.gov。已存档从2022年10月10日开始。已检索10月10日2022年。
\item Stephenson FR (2000年11月)。“亚里士多德观测到的火星月球掩星”。天文学史杂志。31(4):342-344.Bibcode:2000JHA....31..342S。doi:10.1177/002182860003100405。ISSN0021-8286。S2CID125518456。  
\item 哈兰DM (2007年)。卡西尼在土星: 惠更斯的结果。斯普林格。第1页。ISBN 978-0-387-26129-4。
\item McCluskey SC (1998),中世纪早期欧洲的天文学和文化,剑桥大学出版社，pp. 20-21,ISBN 978-0-521-77852-7
\item 李约瑟J,罗南·CA (1985)。中国较短的科学与文明: 李约瑟原文的删减。第二卷 (第三版)。剑桥大学出版社，第187页。ISBN 978-0-521-31536-4。
\item de Groot JJ (1912)。“风水”。中国的宗教-大学主义: 道教和儒学研究的关键。美国宗教史讲座，第10卷。P。普特南的儿子。第300页。OCLC491180。已检索1月5日2016。 
\item Crump T (1992)。日本的数字游戏: 现代日本对数字的使用和理解。日产研究所/Routledge日本研究系列。Routledge. pp. 39-40.ISBN 978-0-415-05609-0。
\item Hulbert HB (1909) [1906]。韩国的通过。Doubleday, Page & Company. p. 426。OCLC26986808。 
\item 库珀K (2023年12月22日)。约翰内斯·开普勒: 解开行星运动的秘密。空间。已检索4月16日2025。




\end{enumerate}